\chapter{荒域}

键盘敲下最后一个句点，叶牧谦长长地吁出一口气，身体向后瘫进椅子里，发出不堪重负的吱呀声。

屏幕上，文档的标题赫然是《问道仙途·基础地理设定综述》。洋洋洒洒数万字，从天界仙城的缥缈云海，到下界十万群山的森然诡谲，山脉走向、江河分布、秘境遗迹……

他把自己能想到的所有关于一个玄幻世界的宏大地理框架都塞了进去，细节丰富，逻辑自洽，甚至随手画了几幅示意地图。

然而，也仅此而已。

社会形态——空白。

势力划分——空白。

修炼体系——空白。

就连网文最重要的故事大纲——还是空白。

叶牧谦揉了揉布满血丝的眼睛，对着屏幕苦笑。

作为一名中文系大三学生，写网文是他最重要的外快来源。不得不说，叶牧谦确实是一个有野心的作者——不，应该说是“设定集写手”，其写作功力实际上并不能完全驾驭这样的书，其“硬塞”感总是过于明显。

上一本书刚刚完结，成绩嘛，不好不坏，总归是平稳落地。编辑催得紧，新书必须尽快跟上，他熬了两晚，总算把最耗时但也最安全的地理设定堆完了。

至于那些需要真正创意和严密逻辑的修炼等级、功法原理、社会矛盾……他脑子里只有一堆从各路小说里扒拉来的、彼此冲突的模糊概念，根本串不成一条能自圆其说的“道途”。

他瞥了一眼窗外泛起的鱼肚白，头痛欲裂。身体和精神都透支到了极限，太阳穴突突地跳，眼前的屏幕开始重影、晃动。

他勉强站起身，想去倒杯水，再睡一觉，却感觉一阵天旋地转。

坏了，低血糖还是猝死前兆？

这个念头刚冒出来，眼前便猛地一黑。不是闭眼的那种黑，而是所有光线、声音、感知被瞬间抽离，坠入无边虚无的绝对黑暗。他甚至没来得及感到恐惧，意识便彻底沉沦。

……

率先恢复的，是嗅觉。

一股清冽到近乎冷冽的空气涌入鼻腔，带着泥土、青草、松针，还有一种……难以言喻的、仿佛被山泉洗过千百年的洁净石头的气息。

这绝不是他那间充斥着外卖盒和灰尘味道的出租屋该有的空气。

紧接着，是触觉。身下硬邦邦的，硌得慌，像是直接躺在粗糙的石板上。微风拂过皮肤，带来明显的凉意。

最后，是听觉。万籁俱寂，但这种寂静并非无声，而是属于深山旷野的、充满细微生命律动的静。远处有隐约的流水潺潺，近处有风穿过林木的飒飒声，偶尔夹杂一两声清脆悠远的鸟鸣。

蝉噪林逾静，鸟鸣山更幽。

这是他的第一反应。

他的第二反应是，他的出租屋环境好像没有这么好。

叶牧谦猛地睁开眼，映入眼帘的是湛蓝到令人心慌的天空，几缕薄云丝絮般挂在天边。

视线下移，是陡峭嶙峋的灰白色山岩，以及岩缝中顽强探出的、姿态奇崛的老松。他正仰面躺在一片天然形成的、不大的岩石平台上。

“我……在哪？”他撑起身体，环顾四周，心脏骤然缩紧。

这里是一座高山的峰顶附近，地势相对平缓的这一小片区域，竟矗立着一座院落。

院墙是就地取材的粗糙山石垒砌，不高，却自有一股历经风雨的沧桑感。院门是简单的木头，虚掩着。

院内，可见两三间同样石木结构的屋舍，瓦片陈旧，檐角却以一种奇异的弧度微微上翘，在晨光中勾出清冷的剪影。

整个院子简陋至极，没有任何奢华装饰，但不知为何，静静立在这云雾缭绕的孤峰之上，竟隐隐透出一股超然物外、不染尘埃的意味。

仙风道骨。

叶牧谦脑子里瞬间蹦出这个词，随即感到一阵荒谬的寒意。这不是他笔下任何一个场景，也绝不属于他认知中的任何现实地点。

穿越？猝死后的幻觉？某种整蛊综艺？

纷乱的念头还没理清，一阵窸窸窣窣的脚步声，伴随着刻意压低的对话声，由远及近，停在了院门外。

“哥，回去吧！这都第几次了？绝念峰，鸟不拉屎的地方，哪有什么仙人遗迹！”一个清脆的女声，带着无奈和疲惫。

“千羽，你不懂。”一个年轻男子的声音响起，低沉而执着，甚至有些魔怔般的激动。

“古籍残卷有载，‘仙缘渺渺，或在群山之末，心诚则灵’。我遍访名山大川不得，或许正是机缘未到，或许……正是要在这无人问津的穷僻之地，方能得见一线天机！

“若无仙缘，不能问道长生，碌碌此生，与草木同朽，有何意义？！”

“长生长生，你就知道长生！那些都是话本传说，骗人的！世上哪有什么飞天遁地、长生不老的仙人？”女子劝道，语气焦急。

“话本？若真是虚妄，为何历代皆有寻觅者？为何总有零星记载流传？我不信！今日既到此地，必要看个究竟……嗯？”

男子的声音戛然而止，似乎极为惊讶，“这里……这里何时竟有了一座院子？上次我来时，分明只有乱石杂草！”

叶牧谦浑身汗毛倒竖。对话内容信息量太大——“十万群山”、“大周国”、“仙人”、“长生”……每一个词都在锤击他脆弱的认知。

而更关键的是，他们发现这个小院了！自己这个穿着现代短袖短裤、来历不明的家伙，就在院子里！

他下意识想躲，可这光秃秃的平台，除了几间屋子无处可藏。就在他惊慌失措时，“咚咚咚”，敲门声响起，格外清晰，敲在他的心尖上。

逃不掉，只能面对。

叶牧谦深吸一口气，努力让自己镇定下来，走到院门前，迟疑了一下，拉开了那扇简陋的木门。

门外，站着一对青年男女。

男子约莫二十出头，身材颀长，穿着一身料子不错的青色劲装，但此刻风尘仆仆，袖口衣角沾着草屑泥点。

他面容俊朗，剑眉星目，只是那双眼睛此刻布满了血丝，闪烁着一种近乎狂热的执着光芒，紧紧盯着叶牧谦，仿佛要将他从里到外看个通透。

这家伙身体瘦弱，从头到脚都不像是个习武之人，唯一像的地方是佩着一把剑——但以叶牧谦粗浅的认知来看，这把剑更像是装饰品。

女子年纪稍小，约莫十七八岁，穿着鹅黄色的裙装，针脚细密而繁复。容貌清丽，眉眼间与男子并无太多相似之处。

她脸上带着长途跋涉的倦色，更多的则是担忧和疑虑，目光在叶牧谦身上飞快扫过，尤其是在他那身格格不入的“奇装异服”上停留片刻，柳眉微微蹙起。

门开的瞬间，空气仿佛凝固了。

青衣男子白言冰的目光，从叶牧谦的短发、黑框眼镜——幸好穿越时还戴着，印花T恤、沙滩裤，一路看到脚上的拖鞋，瞳孔骤然收缩，脸上的表情从惊愕迅速转化为一种难以形容的震撼与狂喜！

“仙……仙人！果然是仙人！”白言冰声音颤抖，脱口而出。在他有限的认知里，这等迥异于凡俗的装扮，出现在这人迹罕至的孤峰秘院，除了传说中的世外高人、游戏风尘的仙人，还能有别的解释吗？

这奇怪的衣装，定是仙家独特的“法衣”！

话音未落，他竟毫不犹豫，推金山倒玉柱般，“噗通”一声直接跪倒在坚硬的石地上，以头触地，纳头便拜：“晚辈白言冰，痴求仙道多年，今日得见仙颜，实乃三生有幸！

“恳请仙师垂怜，收晚辈为徒，传我长生之法，引我踏上问道之途！弟子愿执鞭坠镫，侍奉左右，万死不辞！”

咚咚咚！连着三个响头，磕得实实在在，石板上都见了些许灰印。

叶牧谦被这突如其来的大礼吓得魂飞魄散，差点原地跳起来。

“别！快起来！我哪里是什么仙人！你认错人了！”他慌忙侧身避开，伸手想去扶，又觉得不妥，手僵在半空，语无伦次。

“哥！你疯了！”白千羽也被兄长的举动惊呆了，急忙去拉他，“你快起来！你看他……他这模样，哪里像仙人了？说不定是哪里来的山野怪客，你快起来！”

白言冰却如磐石般跪地不动，反而甩开妹妹的手，抬头仰视叶牧谦，目光炽烈如火：“仙师不必考验弟子！弟子心诚，天地可鉴！此峰名为‘绝念’，向来荒芜。

“此院昨日尚且无存，今日仙师便在此现身，不是仙家手段，又是如何？仙师这身装扮……更是超凡脱俗，非尘世所有！求仙师开恩，收下弟子吧！”说着，又要磕头。

叶牧谦心里叫苦不迭。院子怎么来的他根本不知道！这身打扮在对方眼里成了仙家证明，更是荒谬绝伦。

他连连摆手，急得额头冒汗：“我真不是！我……我就是个普通人，迷路了！不会修炼，更不懂什么长生！你快起来，咱们好好说……”

“仙师若不肯收留，”白言冰的眼神忽然变得决绝而惨烈，他猛地从腰间拔出一柄匕首，寒光一闪，竟直接抵在了自己的心口。

“弟子问道之心已绝，长生无望，苟活于世亦是行尸走肉！不如今日便血溅于此，以明心志！”

“白言冰！你干什么！”白千羽尖叫一声，脸色煞白，想夺匕首却又不敢妄动。

叶牧谦脑子“嗡”的一声，彻底懵了。写网文时，他笔下主角遇到过各种拜师、考验、忠贞不二的桥段；

但当这一幕活生生、血淋淋地发生在眼前，一个活人真的以性命相逼时，那种冲击力完全不是文字可以形容的。

他能看到白言冰眼中那不容置疑的疯狂与绝望，那匕首尖已经刺破了衣衫，再进一步，便是惨剧。

救人！不管怎样，先阻止他！

“住手！我答应你！你把刀放下！”叶牧谦几乎是吼出来的，声音都变了调。

白言冰动作一顿，死死盯着他：“仙师此言当真？不是敷衍弟子？”

“当真！当真！你快放下！”叶牧谦心脏狂跳。

白言冰凝视他片刻，似乎确认了他眼中的惊惶与急切不似作伪，这才缓缓放下匕首，但并未收起，仍握在手中，再次重重叩首：“弟子白言冰，拜见师尊！多谢师尊收录之恩！”

叶牧谦腿都有些发软，靠着门框才站稳。他看着跪伏在地、肩头微微颤抖，不知是激动还是后怕的白言冰，又看了看旁边脸色变幻不定、眼神复杂地盯着自己的白千羽，只觉得一口老血堵在胸口，吐不出来也咽不下去。

这叫个什么事儿啊！

好说歹说，总算让白言冰站了起来，三人进了那间最大的、应该是堂屋的石屋。屋内陈设简单得令人发指：一张石桌，几个石凳，一个空荡荡的土炕，墙角堆着些干柴，除此之外，别无长物，倒是颇为干净。

叶牧谦坐在主位的石凳上，感觉屁股冰凉，心里更凉。白言冰垂手恭立在下首，目光灼灼，满是期待。白千羽则坐在侧面，一双妙目在叶牧谦和兄长之间来回打量，警惕之色未消。

“你们……先说说，这是什么地方？你们又是何人？”叶牧谦强迫自己冷静，开始套取信息。

白言冰恭敬回答：“回禀师尊，此地乃是大周国西境边陲，十万群山山脉最东端的一座山峰，本地人称之为‘绝念峰’。弟子白言冰，这是舍妹白千羽。”他指了指旁边的女子，“我二人……呃，家中经营些生意，略有资财。”

他说这话时眼神闪烁不定，白千羽也偷偷拉他的衣袖。不过叶牧谦本人并没有时间管这个。

白言冰的声音随即变得坚定起来：“弟子自幼慕道，遍寻仙踪而不得，今日幸遇师尊，实乃天意！”

大周国？十万群山？绝念峰？富商（大概）子弟？

当听到“舍妹”的时候，叶牧谦却感觉这对兄妹根本不像。

叶牧谦随口问：“那你们两个年轻人，家里就随便放任你们乱跑？”

白言冰一怔，随即苦笑："家中……管得紧，我们是偷跑出来的。"

白千羽低头整理药囊，轻声补充："兄长执念太深，若不遂他心愿，只怕早晚出事。"

叶牧谦叹了一口气。

得，还是个犟种。

在他看来，谁家丢了两个孩子不着急？这两位倒像是有恃无恐，是根本就没人敢管？

不过眼下还不是研究他们家事的时候。

快速消化完这些信息，叶牧谦又问：“你说寻仙多年，那这世上，可有仙人确切存在的记载？可有修行宗门、长生世家？”

白言冰闻言，脸上闪过一丝黯然，摇了摇头：“不瞒师尊，弟子查阅无数典籍野史，关于仙人的记载大多语焉不详，近乎传说。未曾听闻有确切的修仙宗门或长生世家存世。世人多以之为虚妄怪谈，或方士骗术。”他顿了顿，眼神重新变得炽热，“但弟子坚信其有！今日得见师尊，便是明证！”

叶牧谦心里顿时五味杂陈。

搞了半天，这个世界可能根本就不是个修仙世界！或者说，修仙只是遥不可及的传说，并未形成真正的文明体系！自己这个写书的家伙，居然被一个修仙传说中的狂热信徒，当成了真仙给拜了！

荒唐！滑稽！但又让人笑不出来。

他看着白言冰那充满信仰的眼神，知道任何解释或者澄清在此刻都是苍白的。

对方根本不会信，弄不好又要以死明志。

必须脱身！越快越好！

一个计划，在他脑中迅速成形。

既然你说没有仙人记载，没有修炼体系，那我把我前世印象里那套只有空架子、完全没有具体修炼方法的“问道途”概念抛出来，让你去练。

在一个没有“灵气”——先不管这里到底有没有了，暂时先假设这里没有——或者根本不适合修炼的世界，你注定失败！

到时候，我就可以顺理成章地以“资质不堪，仙缘浅薄”为由，把你打发走，然后我自己再想办法搞清楚状况，寻找回去的路或者在这个世界的生存之道。

对，就这么办！叶牧谦暗自给自己打气。他清了清嗓子，脸上努力挤出一丝属于“世外高人”的淡漠与疏离。

“嗯……本座……”他差点被自己的称呼呛到。

作为在新中国经受过十四年教育的学生，这种话说出来还是太奇怪了！

连忙调整代词，“我云游至此，暂居此院。你既心诚，我便予你一个机会。”

白言冰闻言，激动得浑身一颤，又要下拜，被叶牧谦用眼神制止。

“然，仙道艰难，首重根基与缘法。”叶牧谦开始发挥他编设定的特长，语气变得缥缈起来，“我这一脉，所修之道，名曰‘问道途’。此途重己身，炼体魄，凝精神，悟天地，而非汲汲于外物灵气。”

先把“可能没有灵气”这个坑避开。

“今日，我便传你‘问道途’初基第一境——‘外劲’之法。”

叶牧谦搜肠刮肚，回忆着看过的无数武侠、玄幻小说里关于打基础、练体魄的模糊描述，结合一点粗浅的健身和呼吸知识，开始现场胡编：

“外劲者，由外而内，锤炼筋骨皮膜，激发人身本源之力。其法门，在于特定的呼吸节奏，配合独特的肢体拉伸、负重击打之法，感受肌肉筋膜之震颤，气血运行之鼓荡……”

他说得云山雾罩，夹杂着大量“似松非松”“将紧未紧”“意随身动，气随意走”之类的虚词，根本没有具体可行的步骤。

最后，似乎是为了能“自圆其说”，还故意设置了一个听起来就很玄乎的门槛：“初习此法，需于每日晨曦初露、紫气东来之时，面向东方，以我所述之呼吸法，感应周身三尺内‘微芒’之动。何时能清晰感知并引动一丝‘微芒’淬炼指端，方算入门。此过程，短则数日，长则……数年数十载亦未必可成，全看个人资质与造化。”

在叶牧谦看来，这纯粹是扯淡。“微芒”是什么？他自己都不知道。感知并引动？在一个非修仙世界，这根本就是不可能完成的任务。白言冰注定会一无所获，很快就会知难而退，或者因为“资质太差”被自己劝退。

白言冰倒是听得如痴如醉，虽然许多词汇闻所未闻，但越是艰深晦涩，在他听来越是觉得高深莫测，符合心目中仙家妙法的形象。他根本不去思考可行性，只觉得眼前打开了一扇全新的大门。

“弟子谨遵师尊教诲！定勤加修习，不负师尊厚望！”白言冰躬身应道，脸上洋溢着朝圣般的光彩。

“哥！”白千羽终于忍不住了，站起身，语气带着嗔怪和深深的忧虑，“这……这听起来太过玄虚！你……”

“千羽！”白言冰打断她，神色严肃，“此乃仙缘，岂容置疑？我意已决，你休要多言。师尊肯传授无上妙法，已是天大的恩典！”

他转向叶牧谦，再次行礼：“师尊，弟子这便去寻一僻静处，尝试感应‘微芒’！”

叶牧谦巴不得他赶紧走，自己好静静，便故作高深地点点头：“嗯，去吧。心要静，意要专。”

白言冰兴冲冲地出了院子，真的就在不远处一块凸出的岩石上盘坐下来，面朝东方，闭上眼睛，按照叶牧谦那套颠三倒四的呼吸描述，开始尝试。

堂屋内，只剩下叶牧谦和白千羽。

气氛有些尴尬。

白千羽看着叶牧谦，眼神里的怀疑几乎要溢出来，但她似乎也明白兄长此刻的状态劝不动，只好叹了口气，语气稍微缓和，却依旧带着审视：“这位……先生。我兄长他痴迷此道已久，行事难免偏激，方才多有冒犯，还望海涵。”

“无妨。”叶牧谦摆摆手，心里想着怎么把这姑娘也打发走。他只觉得这姑娘在心情不太激动的时候，措辞有些过分文雅了：谁没事会说“兄长”？

白千羽默默收拾出一间侧屋，又从随身携带的包裹里取出些干粮清水，分给叶牧谦一份，态度谈不上恭敬，却也维持着基本的礼节，只是那双眼睛里的审视意味从未消退。

她动作间，袖口滑落，露出一截皓腕。其上系着一根极细的金链，链坠是一枚小巧的凤纹玉扣，雕工精绝，非寻常匠作可及。

叶牧谦已经一天没吃东西了，正饿着呢。见到白千羽递给他一块大饼，连忙接过；但又觉得狼吞虎咽直接开吃有碍观瞻，于是只能故作矜持，小口吞咽起来。

白千羽见此，心中的怀疑与鄙夷更甚，不动声色地把属于白言冰的那份干粮藏了起来。

两人吃完，再次各个坐定。

“只是，”白千羽话锋一转，目光锐利，“先生所言修炼之法，闻所未闻。千羽自幼随家中医师辨识药材，略通医理，深知人体奥妙，亦有极限。先生之法，似乎……全然不同于任何强身健体或导引之术。”

叶牧谦心里一紧，这姑娘不好糊弄！

他面上不动声色，淡淡道：“仙凡有别，你所知医理，不过是凡尘之见。岂不闻‘夏虫不可语冰’？”但实际上内心疯狂紧张：你怀疑，我当然比你更怀疑我自己啊；究竟是不是编的，我可比你清楚多了！

白千羽被噎了一下，抿了抿唇，显然不服，但也没有再直接反驳。

她沉默片刻，道：“兄长既已拜师，执意留下修炼，我这做妹妹的，也不能独自离去。从今日起，千羽便也留在此处。”

她抬眼看向叶牧谦，目光清亮，“一来，照顾兄长饮食起居；二来……”她顿了顿，语气平静却坚定，“千羽也想亲眼看看，先生所传的‘仙法’，究竟有何神异之处。想必先生，不会介意吧？”

叶牧谦听出了她的弦外之音：照顾兄长是假，留下来监视自己这个“来历不明的怪人”才是真。他心中苦笑，却也无法拒绝。难道说“你赶紧走，别妨碍我忽悠你哥”？

“随你。”他只能故作淡然，心里却在哀嚎：这下好了，走不脱了，还多了个监工！

只能指望白言冰尽快修炼失败，自己才能逃跑！

他走到门口，望向远处岩石上那个如雕塑般一动不动、认真尝试感应所谓“微芒”的青色身影，又瞥了一眼屋内开始默默收拾、目光却不时瞥向自己的鹅黄衣裙少女，心中一片茫然。

又憋了半日，白千羽终究没忍住，开口问道：“先生所说的‘微芒’，究竟是何物？是天地间的一种精气，还是人体内潜藏的一点灵光？”

她自幼接触药材医理，对人体气血、自然物性有着朴素的认知，叶牧谦那套玄乎的说法，既让她怀疑，也勾起了她探究的本能。

叶牧谦心里一虚，面上却绷着：“此乃道之初始，不可言，不可名。说似一物即不中。”他干脆甩了句禅宗话头，企图蒙混过关。

白千羽若有所思，却没有继续追问，转而道：“兄长自幼体弱，家中曾请名医调理，言其先天心脉有亏，不宜剧烈运动，更忌忧思过度。他这般痴迷……耗神枯坐，于身体可有妨碍？”

这问题更具体了，叶牧谦头皮发麻，他懂什么“心脉有亏”？

但对方看着自己，等一个回答。

他只能搜刮前世看过的那些中医养生文章、武侠小说里的经脉理论，含糊道：“仙道修炼，本就是逆天改命，弥补先天不足。

“外劲初基，看似锤炼体魄，实则首要在于‘意’与‘息’的调和。呼吸深长匀细，意念专注守一，本身便有安神定志、导引气血之效。至于具体关隘……因人而异，需静观其变。”

他一边说，一边心里打鼓。

这纯粹是车轱辘话。

然而，白千羽听着，眉头却微微舒展开一些，低声自语：“意念专注，呼吸深长……导引气血……似乎与《素问》中‘恬淡虚无，真气从之’有些相通之处，只是更为主动……”她看向叶牧谦的目光，少了几分尖锐，多了些思索。

叶牧谦见状，暗暗松了口气，同时也有些荒谬感：这都能蒙过去？

第一天在一种诡异平静中度过。白言冰除了喝水进食，几乎全天都在那块岩石上打坐。叶牧谦则度日如年，期盼着对方尽快失败放弃。

然而，第二天清晨，天刚蒙蒙亮，叶牧谦就被屋外一阵压抑着兴奋的低呼惊醒。

他披衣出门，只见白言冰站在院中，对着昨夜被他拿来练习“感受气血鼓荡”而拍击了无数次的一块青灰色山岩，缓缓吸气，然后吐气开声，一掌按了上去。

动作并不快，也没有影视剧里那种光影特效。

但就在他手掌与岩石接触的刹那，叶牧谦似乎听到了一声极其细微的、仿佛硬物被挤压的“嘎吱”声。白言冰收回手，只见那原本光滑的岩面上，赫然留下了一个浅浅的、但清晰可见的掌印，边缘甚至有些许石粉簌簌落下！

院子里一片寂静。

白千羽刚刚打水回来，看到这一幕，手中的木桶“哐当”一声掉在地上，清水洒了一地。她瞪大眼睛，看看岩石上的掌印，又看看兄长微微颤抖、却明显洋溢着狂喜的手掌，最后难以置信地望向叶牧谦，红唇微张，一个字也说不出来。

叶牧谦更是如遭雷击，僵在原地，脑子一片空白。

成功了？

徒手在石头上留下印子？

这他妈是什么怪力？

我编的那些玩意儿……真的能练？

这世界到底怎么回事？

不是说没有仙人记载吗？

这不符合物理定律啊！

他快步上前，捡起那块石头查看，却发现这石头可太石头了；想以常人体力在上面留下痕迹，那无异于痴人说梦。

“师……师尊！”白言冰转过身，脸上因为激动而涨红，眼眶甚至有些湿润。

他快步走到叶牧谦面前，再次跪下，声音哽咽：“弟子……弟子似乎感受到了！昨日午后，弟子依照师尊所言调整呼吸，意念存想周身，初时只觉得烦躁困顿；但入夜之后，心神渐宁，忽然觉得指尖似乎有微弱暖流划过，似有似无。

“我便牢记那感觉，继续存想。就在刚才，晨曦照身的那一刻，那暖流骤然清晰了一丝！我福至心灵，顺势将其导向手掌……便、便有了这般效果！师尊，这便是‘微芒’吗？弟子……弟子入门了吗？”

叶牧谦看着白言冰那纯粹而炽烈的眼神，听着他详细的描述，只觉得一股寒气从脚底板直冲天灵盖。他描述的感受，分明就是许多玄幻小说里“气感初生”的桥段！

可那是小说啊！

是胡诌的啊！

“呃……嗯，不错。”叶牧谦喉咙发干，勉强挤出几个字。

他感觉自己必须说点什么：“初窥门径，算是……摸到了一点外劲的边儿。但切不可自满，此乃万丈高楼第一块砖，距离真正掌握‘外劲’，收发由心，还差得远。”他本能地开始泼冷水，拖延节奏。

“弟子明白！定当时刻谨记师尊教诲，勤修不辍！”白言冰重重叩首，喜悦之情溢于言表。

白千羽此时也回过神来。她先抓起兄长的右手查看，掌心只有些许红印，连皮都没破；又走到岩石边，仔细摸了摸那个掌印，确信那是她兄长刚刚留下来的。

她抬头看向叶牧谦，眼神彻底变了，之前的怀疑被巨大的惊诧和一丝隐隐的敬畏取代。

能让人在一天一夜之内发生如此脱胎换骨般变化的“法门”，已经超出了她所能理解的医学范畴。

“先生……真乃神人。”她涩声说道，这一次，语气带上了明显的敬意。

叶牧谦却一点也高兴不起来，只觉得事情彻底失控了。他胡乱点了点头，转身回了屋子，需要静静。

然而，刚一坐下，他自己也察觉到一丝异样。

似乎……身体轻松了一些？

昨天熬夜穿越、心惊胆战带来的沉重疲惫感，消散了大半。不仅如此，他下意识握了握拳，感觉力量好像也大了一点？视线扫过屋内，似乎比往常更清晰些，耳朵也能听到更远处山风吹过松针的细微声响。

是心理作用？还是……

一个让他毛骨悚然的念头浮现：难道白言冰修炼有成，我自己也会跟着变强？

“道行反哺”？

可这反哺的依据是什么？就因为我胡编了一套东西，还被他练成了？

为了验证，他偷偷尝试了一下。集中精神，对着桌上的一个石碗，低声念叨：“飞起来。”

石碗纹丝不动。

挠了挠头，寻思可能是这东西太重了，改用木头。

依然失败。

他又尝试回忆更多前世看过的修炼口诀，什么“气沉丹田”“周天运转”，默念半天，体内毫无感应，也没有任何力量涌现。

“看来不是‘言出法随’……”叶牧谦松了口气，但更多的疑惑涌上心头。

不是心想事成，那这身体的变化从何而来？难道真的只跟白言冰的修炼进度挂钩？

这个发现让他坐立不安。原先的计划彻底破产了。白言冰不仅练成了，还展现出了超常的“天赋”，而自己这个“师尊”，竟然也跟着沾光？这让他还怎么理直气壮地用“资质太差”赶人走？

接下来的几天，叶牧谦陷入了更深的焦虑。白言冰进展神速，那所谓的“外劲”运用越来越纯熟，已经能比较稳定地在石头上留下痕迹，甚至跳跃、奔跑、舞剑的身手明显矫健了许多。

而叶牧谦自己也清晰地感觉到，力量、速度、耐力，乃至精力，都在稳步提升，虽然幅度远不如白言冰夸张，但这确确实实是超越他原有亚健康大学生体质的增强。

叶牧谦知道，不能再这样下去了。白言冰的外劲眼看就要“圆满”——虽然他自己都不知道圆满标准是啥。

他必须抛出新的东西，继续拖延这古怪的“徒弟”。

\vspace{2em}

第七天，白言冰兴奋地向妹妹展示，自己已经能一掌拍碎一块脸盆大的石块。

叶牧谦心知，时机已经成熟，于是摆出了一副严肃的表情：

“外劲粗通，不过是打熬气力的法门，算不得真正踏入‘问道途’。”

白言冰的脸顿时垮了下来，不过立刻被更大的求知欲取代。

“请师尊指示。”

叶牧谦也没想到他会如此快地接受了这个现实。不过这样也好，可以直接把预先打好的腹稿拿出来用了：

“外劲之后，便是‘内劲’：需由外而内，淬炼根骨，滋生内劲。此乃筑基之本，关乎日后道途远近，至关重要。”

“根骨？内劲？”白言冰眼睛发亮，如同听到福音。

“然也。”叶牧谦负手而立，望着远山云雾，一副高人风范。

“人之资质，显于根骨。先天所定，后天也可打熬。根骨分四品：凡骨、灵骨、魂骨、天骨。凡骨者，众生芸芸，修炼事倍功半；灵骨者，百里挑一，可窥道途门径；魂骨者，万中无一，有望登堂入室；天骨者……传说而已，得之可直指大道本源。”

他故意把等级设得多，标准设得玄，尤其是最后的“天骨”，直接定义为传说，就是为了埋下一个几乎不可能达成的目标，用来无限期拖延。

“淬炼根骨，便是以内息温养骨骼、骨髓，涤除杂质，激发潜能。这是一个水磨工夫，急不得。何时你能内视自身，感知骨骼如玉般温润，隐有光华，便是踏入灵骨之境。届时，内劲自生，力量圆转，远非蛮力可比。”

叶牧谦说得天花乱坠，心里却想：内视？感知骨骼如玉？这比感应“微芒”还玄乎！

你就慢慢琢磨去吧，没个十年八年，休想有什么感知。

然而，他再次低估了白言冰的决心，也低估了这个世界的诡异。

白言冰没有任何质疑，将叶牧谦关于“内息温养”、“意守骨骼”、“呼吸节奏配合气血搬运”等一堆东拼西凑、漏洞百出的“法诀”奉若至宝，开始了新一轮的苦修。这一次，他不再需要面朝东方，但修炼更加枯燥，大部分时间都是静坐冥想，配合缓慢而奇特的肢体拉伸与呼吸。

叶牧谦提心吊胆地观察着。

五天过去，白言冰除了精神更显凝练，并无特异表现。叶牧谦稍安。

第十天，白言冰在静坐时，周身隐隐有热气蒸腾，持续时间不长，但被细心观察的叶牧谦和白千羽注意到了。白千羽惊异，叶牧谦心惊。

第十五天，清晨。白言冰从彻夜静坐中醒来，缓缓睁开双眼。那一瞬，叶牧谦仿佛看到他眼中有一抹极淡的、温润的光泽闪过，随即隐去。

白言冰站起身，活动了一下筋骨，全身发出一连串清脆而和谐的“噼啪”声，如同玉珠落盘。他随意一拳挥向空中，竟带起了清晰的破风声，而且动作流畅自然，仿佛力量是从骨骼深处涌出，贯通全身。

“师尊，”白言冰走到叶牧谦面前，语气带着压抑的激动和一丝不确定，“弟子昨夜静坐，恍惚间似能‘看’到自身手臂骨骼，其色似有微光，温润不刺眼。今日发力，亦觉与往日不同，劲力似乎更整了。不知这算不算是……”

叶牧谦看着他清澈又期待的眼神，听着那描述，感受着白言冰身上确实有所不同、更显凝实的气息，一颗心直往下沉。

灵骨？！半个月？这就灵骨了？

我设定的“百里挑一”、“水磨工夫”呢？

这世界的时间流速和修炼难度，是在临期食品仓储店打折促销的吗？

与此同时，就在白言冰说出“似有微光，温润不刺眼”的瞬间，叶牧谦自己体内，也蓦地传来一阵奇异的温热感，仿佛有什么东西在骨骼深处被唤醒、熨帖而过。

紧接着，他明显感觉到自己的力量又涨了一截，五感似乎也敏锐了一丝，昨天搬运柴火时不小心蹭到石头的隐隐痛感，此刻也消失无踪。

反哺！又来了！而且这次感觉更明显！

叶牧谦心中警铃大作。一次是巧合，两次就绝不是了！自己身体的变化，绝对与白言冰的修炼突破直接相关！自己胡编乱造的境界，被他练成了，然后自己这个“创始人”就能分到好处？这是什么原理？

巨大的困惑和隐隐的恐惧攫住了他。

但面对白言冰殷切的目光，他只能压下所有情绪，面无表情地点点头：“嗯，初具灵骨之象，算是入门了。内劲已生，日后需勤加运转温养，使其茁壮。”

白言冰大喜过望，又要行礼，被叶牧谦止住。

“不过，”叶牧谦话锋一转，开始铺设障碍，“灵骨之境，亦有深浅。需得将周身二百零六块骨骼，尽数淬炼至同等程度，方算灵骨圆满，方可考虑下一品‘魂骨’的淬炼。此过程，需极耐心，戒骄戒躁。”他直接把淬炼范围扩大到全身骨骼，难度指数级上升。

白言冰毫无畏难之色，肃然应道：“弟子谨记！”

然而，正如叶牧谦所期望——也是他所害怕——的那样，在初步踏入“灵骨”，感受到内劲滋生的神奇之后，白言冰的进展陡然慢了下来。

全身骨骼的淬炼，远比局部感知要困难、繁琐得多。他陷入了漫长的闭门造车期。尽管日夜苦修不辍，但那种明显的、阶段性的突破感再也没有出现。他变得有些焦躁，静坐时眉头时常紧锁。

叶牧谦一边暗中观察着：反哺似乎也进入了平缓期；自己身体虽持续缓慢增强，但再无那次明显的温热感；一边也在焦急地思考下一步。白言冰卡住了，这是好事，可以拖延时间。但总不能一直卡在这里，万一这小子钻牛角尖，或者觉得是自己教得不好……

机会很快来了。

白千羽每隔几日会下山，去距离绝念峰三十余里外的一个小镇采买些必需品，也顺便打听些消息。她用的却不是叶牧谦以为的碎银或铜钱，而是一种印有复杂水印的银票；随着必需品带来的，甚至还有少量人参等名贵药材，她对此的解释是“家里商队恰好路过”。

但这一次她回来的时候面带忧色。

“兄长，叶先生，”她语气有些急促，“山下不太平。听说最近有一伙山贼流窜到咱们这附近，占了北边黑风岭的老寨，已经劫掠了好几个村庄，伤了不少人。镇上的巡防营去剿过两次，都吃了亏。据说那伙贼人里，有几个身手相当厉害的头目，不是普通兵丁能对付的。”

山贼？高手？叶牧谦心中一动。

他立即叫来了白言冰。

“修炼之道，并非一味枯坐。实战磨砺，生死之间的感悟，往往胜过十年苦修。这伙山贼，倒也算是个不错的历练对象。”

白言冰精神一振。“先生可是希望我去剿匪？”

“然也。不过，”叶牧谦话锋一转，“既是历练，便需独自面对。你一人前去，不得借助任何外力，包括你妹妹。我要你凭借这半月所学，独自解决此事。可能做到？”

“先生！”白千羽失声惊呼，“兄长他虽有些进境，但毕竟修炼日短，那伙贼人凶悍，还有高手……怎能让他独自涉险？”

白言冰却毫不犹豫，斩钉截铁道：“弟子能做到！请师尊放心！”

叶牧谦看着白言冰坚定的眼神，强行硬起心肠。

“好。那便如此定了。明日凌晨出发，给你三日时间。三日之后，无论成败，我要在此处见到你。”

他顿了顿，感觉良心在痛，于是补充道，“若事不可为，保命为先，可退回来再从长计议。”

白言冰只当是师尊的考验与关怀，热血上涌，用力点头：“弟子遵命！定不辱命！”

白千羽还想说什么，却被叶牧谦淡淡的眼神制止了。她看着兄长兴奋又坚毅的侧脸，再看看叶牧谦那看似平静无波、高深莫测的脸，满心忧虑化作一声沉重的叹息。

\vspace{2em}

是夜，叶牧谦辗转难眠。白言冰擦拭着他那柄装饰性大于实用性的佩剑，神情专注。白千羽默默为他准备干粮、水囊和金疮药。

凌晨，天色未明，山间雾气弥漫。白言冰换上了一身便于行动的深色衣服，背好行囊，佩上长剑，对着叶牧谦所在的屋子方向，郑重地行了一礼，然后转身，毫不犹豫地踏入了下山的小径，身影很快被浓雾吞没。

叶牧谦站在窗边，看着他消失的方向，心脏在胸腔里沉重地跳动。那点阴暗的期盼和巨大的负疚感交织折磨着他。

\vspace{2em}

我这么做……是不是太过分了？

\vspace{2em}

他将一个相信自己是仙人、毫无保留修炼自己胡诌功法的青年，推向了一个可能有真实生命危险的境地。而这，只是为了自己那点不可告人的私心：希望他失败，希望他知难而退，甚至希望他再也回不来，一了百了。

“先生。”白千羽不知何时来到门口，声音有些发颤，“兄长他……真的不会有事吗？”

叶牧谦转过身，看着少女苍白的脸和通红的眼眶，强迫自己压下所有情绪，脸上露出一种他自认为应该是“仙师”该有的、洞悉一切的平静。

“道途之上，各有缘法，各有劫数。”他声音低沉，仿佛蕴含着无尽玄机，“此去，是他的劫，也是他的缘。静候便是。”

说完，他不再看白千羽，走到石桌前坐下，闭上眼，仿佛入定。

只有他自己知道，那平静的外表下，是怎样一颗七上八下、备受煎熬的心。

绝念峰上，晨雾缭绕，仿佛也笼罩在了叶牧谦纷乱的心头。道行反哺的奥秘尚未解开，首次由他亲手促成的冲突已然拉开序幕。而他这个冒牌师尊，只能在忐忑与愧疚中，等待着一个他既期盼又害怕的结果。

\vspace{2em}

三日，于叶牧谦而言，漫长得如同三个春秋。

\vspace{2em}

第一日，他在表面的平静与内心的翻江倒海中度过。

白千羽几乎寸步不离地守在山崖边，眺望着兄长消失的方向，茶饭不思，眼眶始终是红的。

叶牧谦则将自己关在石屋内，强迫自己继续“推演”那套虚无缥缈的“问道途”，实则心乱如麻，只是发呆。每一次山风吹动门扉的轻响，都会让他惊跳起来，以为是白言冰归来的脚步声，或是……更坏的消息。

他不断回忆起白言冰跪拜时那炽烈的眼神，练成外劲时的狂喜，以及下山前那毫无保留的信任与坚定。

愧疚如同藤蔓，缠绕心脏，越收越紧。

“我这么做，与那些利用他人信任、草菅人命的恶徒有何区别？”

这个念头反复啃噬着他。可木已成舟，他只能继续扮演那个高深莫测、一切尽在掌握的“仙师”，用沉默和偶尔几句玄乎的“静心”、“等待机缘”来安抚濒临崩溃的白千羽，也麻痹自己。

\vspace{2em}

第二日，焦虑升级为恐惧。

山下并无消息传来，但山间的雾气似乎都带上了一丝不祥的血色（或许只是心理作用）。

白千羽憔悴不堪，终于忍不住带着哭腔质问叶牧谦：“先生，若兄长真有闪失……您真的不担心吗？您传授他法门，难道就只是为了让他去送死历练吗？”

这话如同刀子扎在叶牧谦心上。

他背对着白千羽，望着窗外层峦叠嶂，声音干涩：“道途艰险，各有命数。担心……无用。”

只有他自己知道，袖中紧握的拳头，指甲已深深陷进掌心。

他开始后悔，无比后悔。那点阴暗的期盼早已被巨大的负罪感淹没，他现在只希望那个倔强的青衣青年能全须全尾地回来，哪怕他修炼失败，哪怕他从此不再认自己这个师傅。

\vspace{2em}

第三日，黎明前最黑暗的时刻。

叶牧谦彻夜未眠，白千羽也靠在门边昏昏沉沉。就在天际将亮未亮、山林间万籁俱寂之时，一阵轻微却异常清晰的脚步声，踏碎了凝固的焦虑，由远及近。

那脚步声沉稳，甚至带着一种奇特的韵律，与山贼溃逃或负伤踉跄的慌乱截然不同。

叶牧谦和白千羽几乎同时弹起，冲向院门。

晨雾被一道身影破开。

白言冰回来了。

深色的衣衫沾染了暗沉的血迹、尘土与草汁，多处破损，尤其是左肩和肋下的布料被利刃划开长长的口子，露出里面只有浅浅皮外伤的皮肤。

他脸上带着疲惫，但那双眼睛却亮得惊人，如同被山泉洗过的星辰，锐利、沉静，又隐含着某种经过淬炼后的坚毅。他手中提着那柄长剑，剑身已擦拭过，但靠近剑柄的吞口处仍残留着难以洗净的暗红。

“哥！”白千羽第一个扑上去，抓住他的手臂上下查看，声音颤抖，“你……你没事？伤到哪里了？”

白言冰对妹妹露出一个安抚的笑容，轻轻摇头：“皮外伤，不妨事。”他的目光越过白千羽，投向站在门口、脸色苍白的叶牧谦，快步上前，单膝跪地，抱拳道：“师尊，弟子归来复命。”

叶牧谦喉咙发紧，一时间竟说不出话，只是死死盯着他，确认他真的四肢健全、神志清醒。

“黑风岭的山贼，”白言冰的声音平静无波，却字字清晰，砸在清晨冰冷的空气中，“共计五十三人，包括三名匪首，已尽数伏诛。贼寨已焚。被掳掠的百姓，弟子已让他们各自归家。”

尽数伏诛？五十三人？三名匪首？已焚？

白千羽捂住了嘴，眼睛瞪得滚圆，难以置信地看着兄长，仿佛第一次认识他。那个自幼体弱、需要她时时看顾的兄长，如今却用如此平淡的语气，陈述着一人一剑荡平数十悍匪的骇人事实！

叶牧谦更是如遭雷击，脑子里“嗡”的一声。

全歼？他一个人？这怎么可能？！

就算他有了“灵骨”和内劲，可那是五十多个穷凶极恶、其中还有“高手”的山贼！

不是五十棵白菜！

他预想过白言冰可能负伤逃回，可能狼狈不堪，甚至可能就此折在那里，但唯独没想过是这般轻描淡写的“全歼”！

“你……详细说来。”叶牧谦听到自己的声音有些飘忽，“坐起来说。”

白言冰调整了下姿势，改跪姿为盘腿，略一沉吟，开始叙述。他的语调没有什么起伏，仿佛在说一件与己无关的事，但其中的凶险，却让听者脊背发寒：

“弟子循着千羽所述路径，于第一日午后抵达黑风岭附近。未敢贸然进山，先在山下遭劫的村庄打听。村民言，匪巢位于岭上一处易守难攻的险地，有三道哨卡，匪首三人。

“其中，大当家使一柄九环鬼头大刀，势大力沉；二当家擅使一对短戟，身形灵活，阴险狡诈；三当家则是个使毒镖的侏儒，防不胜防。其余匪众亦多是亡命徒。

“弟子趁夜色掩护摸上山。第一道哨卡两名匪徒，正在打盹。弟子本欲绕过，其中一人忽然惊醒。弟子不及细想，以内劲贯于指端，弹出一粒石子，击中其昏睡穴。另一人闻声欲起，弟子已掠至身前，以掌缘切其颈侧，二人皆昏厥。

“第二道哨卡设在陡坡上，有四人，较为警觉。弟子欲攀岩迂回，不慎踩落一块松石。声响惊动匪徒，他们持刀呼喊围来。弟子知不能留手，拔剑迎上。

“初时有些生疏，只以外劲蛮力格挡劈砍，震得匪徒刀剑脱手。后来忆起师尊所言‘内劲圆转’，尝试将体内那股温热气流导引至手臂。剑招顿时轻灵许多，力道却更透。寻常匪徒，往往一招之间，便被震断兵器，或点中要穴倒地不起。”

叶牧谦听着，心中骇然。

内劲圆转？那只是他随口胡诌的形容词！白言冰竟然真的在实战中摸索出了应用之法！

“突破第二道哨卡后，匪寨已近在眼前。其时天色微明，寨中喧哗，似在聚餐饮酒。弟子潜至寨墙下，正欲翻入，忽觉脑后生风，急侧身避让，一枚泛着蓝光的毒镖擦着耳畔飞过，钉入树干，滋滋作响。

“原来是那使毒镖的侏儒三当家藏在暗处，接连发射毒镖，角度刁钻。弟子挥剑格挡，叮当之声不绝。其余匪众被惊动，蜂拥而出。

“混战之中，匪众悍不畏死，刀枪并举。弟子将内劲灌注双腿，步法变幻，于人群中穿梭，剑光闪处，必有匪徒倒地。但匪众太多，渐渐被围在垓心。

“那使短戟的二当家突然从人缝中窜出，双戟如毒蛇吐信，直刺弟子肋下与后腰。彼时旧力已尽，新力未生，避无可避。只得勉强扭身，以内劲硬抗。

“二当家那戟刺破了弟子衣衫，但触及皮肤的瞬间，弟子骨骼似乎自发地微微一震，一股反震之力涌出，竟让那精铁打造的戟尖难以深入，只留下两道白痕，划破了油皮。

“二当家随即愣住，弟子抓住这电光石火的机会，反手一剑，剑身拍在其手腕上，骨折声响起，短戟落地。紧接着一脚踹出，二当家惨叫着跌入人群。

“大当家见状怒吼，挥舞那柄沉重的鬼头大刀劈头砍来，风声呼啸。弟子不敢硬接，施展身法游斗。刀剑相交数次，火花四溅。弟子觉其力大势沉，但招式略显笨拙。遂卖个破绽，诱其全力一刀劈空，身形踉跄。弟子揉身而上，以内劲灌注剑尖，点向其胸口膻中穴。”

白言冰描述到这里，眼中闪过一丝奇异的光彩，“就在剑尖及体的刹那，弟子福至心灵，体内那股温热内劲以前所未有的速度奔涌，瞬间贯通手臂，透过剑尖……”

他没有详细描述结果，但叶牧谦和白千羽都能想象。那势大力沉、凶名赫赫的大当家，恐怕连惨叫都未能发出，便被这凝聚了初生内劲的一点击溃了生机。

“大当家毙命，余匪胆寒，开始溃逃。弟子追击，那侏儒三当家欲趁乱再施毒镖，被弟子以石子先一步击中手腕，毒镖落地，后被斩首。其余匪众，或死于剑下，或坠崖而亡，无一走脱。”

白言冰的声音低了下去，“随后，弟子寻到被掳掠的妇孺，将他们放出，又点燃贼寨粮草物资，大火烧了半日方熄。待处理完诸事，已是第二日夜间。弟子在山泉边洗净血污，略作调息，便连夜赶回。”

叙述完毕，院子里一片死寂。只有山风穿过松林的呜咽。

一人一剑，一夜一日，连破三关，斩杀五十三悍匪，其中三名头目堪称“高手”，自身仅受微不足道的擦伤……这已不是“碾压”可以形容，这简直是虎入羊群，是境界上的绝对差距！

叶牧谦看着坐在面前、平静陈述这一切的青年，第一次如此清晰地认识到，自己胡编乱造的“问道途”，在这个世界，究竟造就了一个怎样的怪物！

而自己，竟然曾希望这个对自己奉若神明的“怪物”去送死！

巨大的惊诧之后，是更汹涌的羞愧，几乎要将他淹没。

他张了张嘴，想说点什么，却发现任何言语在此刻都苍白无力。

最终，他只是干涩地道：“起来吧。做得……不错。”这称赞，连他自己都觉得虚伪。

白言冰却如获至宝，脸上绽开纯粹的笑容：“多谢师尊！”

殊不知，这一“谢”，反而让叶牧谦的心更痛了几分。

白言冰汇报完屠戮山匪的过程，从怀中取出一块令牌，递给白千羽：“从大当家身上搜出的，你看。”

白千羽接过，脸色骤变：“这是……边军的制式腰牌！怎么会流落到山贼手中？”

白言冰沉声道：“不止。那大当家的刀法，有军中套路的影子。这批山贼，怕是没那么简单。”

叶牧谦倒是警觉起来：这小子，怎么对军伍之事这么门儿清？寻常富商子弟不更应该考虑山贼手里的账本吗？随即开口询问：“你们到底是什么人？”

白言冰犹豫片刻，道：“师尊，弟子……弟子家中与朝中有些关联。弟子不敢欺瞒师尊，但确有难言之隐。待此间事了，弟子必如实相告！”

随即白言冰站起身，忽然身形微不可察地晃了一下，眉头轻轻蹙起。

“哥？”白千羽立刻察觉。

“无妨，只是……有些奇异的感受。”白言冰摸了摸自己的胸口，又活动了一下肩膀，眼中流露出困惑与思索，“与那大当家最后交手时，内劲勃发之后，体内骨骼似乎……有些发痒，又有些清凉，仿佛有什么东西在生长、在蜕变。方才赶路时便有感觉，此刻愈发明显了。”

叶牧谦心中猛地一跳！难道……

“你且回房，静坐内观，细细体会，不可急躁。”他连忙吩咐，心中那点愧疚被新的惊疑取代。

白言冰依言，向二人行礼后，便回了自己那间侧屋，关上了房门。

接下来的大半天，叶牧谦和白千羽都无心他事。白千羽在兄长房外徘徊，时而贴耳细听，里面却寂然无声。叶牧谦则坐立不安，既期待又害怕。期待的是验证“反哺”的猜想，害怕的是白言冰再次突破，将他这个“师尊”架在更高的火上烤。

夜幕降临，星斗满天。

约莫子时前后，白言冰房中突然传出一阵极其轻微、却又清晰无比的“嗡”鸣声，仿佛优质玉器被轻轻敲击后的余韵，悠长而纯净。紧接着，一股难以言喻的气息弥漫开来，并不强烈，却让院中的叶牧谦和白千羽同时感到心神一宁，仿佛被清冽的泉水洗涤过一般。

房门“吱呀”一声打开。

白言冰走了出来。他换了一身干净的青色衣衫，整个人似乎瘦削了一丝，但身姿更为挺拔，如同经过精心打磨的玉竹。肌肤在月光下隐隐流转着一层极其淡薄、温润的光泽，双目开阖间，神光内蕴，深邃宁静。

他站在那里，自然流露出一种沉凝如山、又灵动如水的奇异气质。

“师尊，千羽。”白言冰开口，声音比往日更加清越平和，“弟子……似乎突破了。”

“突破？”白千羽急问，“是内劲大成了吗？”

白言冰摇了摇头，看向叶牧谦，眼中闪烁着领悟的光芒：“弟子静坐内观，那骨骼发痒清凉之感愈烈。意念集中之下，恍惚间见自身二百零六块骨骼，通体呈现一种温润的淡金色，质地紧密，隐有玄奥纹路若隐若现。

“弟子气血流转其间，畅通无阻，生机勃勃。心念一动，内劲自骨骼髓海中滋生，沛然流转，圆融一体，远胜从前十倍！师尊，这莫非就是……‘魂骨’之境？”

魂骨！万中无一，有望登堂入室的魂骨！他半个月前才刚入灵骨，一次生死历练，一夜闭关，竟然直接跨入了魂骨！

叶牧谦只觉得一股寒气从尾椎骨直冲头顶，荒谬感与现实感激烈冲撞，让他几乎站立不稳。他设定的，用来无限期拖延的、万中无一的门槛……就这么被跨过去了？而且是在这种情形下？

与此同时，就在白言冰说出“魂骨”二字，身上那温润沉凝气息稳固下来的瞬间，叶牧谦体内，那股熟悉的温热感再次涌现！但这一次，远比“灵骨”时强烈得多！

仿佛有一股温热的暖流，从他自己的骨骼深处、从每一个细胞中被唤醒、激发出来！四肢百骸发出细微的欢鸣，力量感、轻盈感、五感的敏锐度，都在以清晰可辨的速度提升！

他甚至能听到更远处夜枭归巢的振翅声，能看清黑暗中树叶的细微纹理，轻轻一握拳，指节发出清脆的响声，蕴含的力量让他自己都感到心惊。

反哺！确凿无疑的反哺！而且，白言冰突破的境界越高，自己获得的反哺就越强！

巨大的惊诧之后，那股一直被压抑的羞愧，如同火山般喷发出来，瞬间淹没了叶牧谦。他看着眼前这个气质蜕变、眼中带着对“道”的虔诚与对自己毫无保留信任的青年，想起自己之前那些阴暗的算计、冷漠的推诿、甚至希望他遭遇不测的卑劣念头……

“我……我真不是个东西！”这句话在叶牧谦心中轰然作响，震得他灵魂发颤。

白言冰如此单纯，如此倔强，将自己胡诌的一切奉为圭臬，以性命相托，在生死间砥砺，竟真的闯出了一条前所未有的路。而自己这个“师尊”，除了给予一套自己都不信的框架，还给了什么？怀疑、利用、甚至几乎将他推向死地！

面对这样的徒弟，自己还有什么脸面只想着逃离、敷衍、利用？

叶牧谦缓缓闭上了眼睛，深吸了一口清冷的山间空气。当他再次睁开眼时，眸中的迷茫、焦虑、算计，渐渐沉淀下去，取而代之的是一种前所未有的清明与责任。

是的，责任。

既然来了，既然收了，既然这胡编的“问道途”真的能走通，甚至与自己息息相关……那就，好好培养吧。

能培养到哪一步，就培养到哪一步。

至少，要对得起这份毫无保留的信任，对得起这青年眼中纯净的求道之光。

至少，要问心无愧。

当然，试图回家当然依然是他的主要任务，但至少现在，眼见得是没办法的。

那就暂时放下“回家”或“逃离”的思维吧。

叶牧谦，这个来自异世的灵魂，终于从心底里，接受了“师尊”这个身份，接受了白言冰这个徒弟。

当然，这也是部分因为这片天地似乎确实接纳了他的存在，至于自己的出身——暂时也先别告诉徒弟们吧。

他只是隐隐觉得自己这套逻辑一定能走通。至于为什么这么觉得，他自己也不知道。

他看向白言冰，目光不再有躲闪和虚浮，而是变得沉稳、温和，带着一丝真正的赞赏与关切。

“不错。”叶牧谦的声音平和了许多，也真诚了许多，“初入魂骨，根基未稳。需沉心静气，巩固境界，细细体悟此番突破的所得，尤其是生死之间的感悟，与内劲骨髓相融的玄妙。路还长，切莫自满。”

这不再是敷衍的套话，而是真正的师长叮嘱。

白言冰敏锐地察觉到了师尊语气中那微妙却真实的变化，心中涌起一股暖流，再次郑重行礼：“弟子谨遵师命！”

一旁的白千羽，看着气质迥异的兄长，再看向目光变得温润平和的叶牧谦，心中最后那点戒备和疑虑，也如同阳光下的冰雪，悄然消融了。

兄长真的变强了，强得超乎想象，而这一切，都源于这位看似年轻、行事古怪的“叶先生”传授的法门。他不是骗子，不是怪客，而是一位真正身怀莫测玄功、教导方式特立独行的世外高人！

她走到叶牧谦面前，盈盈下拜，这一次，是真心实意的敬礼：“千羽往日多有疑虑冒犯，请先生海涵。多谢先生传授兄长真法，救命、成全之恩，没齿难忘。”

叶牧谦坦然受了这一礼，虚扶一下：“起来吧。你兄长心志坚毅，是他自己的造化。你这些时日照顾周全，亦是有功。”

三人之间的气氛，第一次如此和谐、坦然，带着一种真正同门师徒间的羁绊感。

白言冰巩固魂骨境界数日，气息日益沉凝。他依然不忘叶牧谦当初划分的“根骨四品”，在稳固境界后，又尝试感应那传说中的“天骨”之境，但毫无头绪，仿佛那真的只是一个虚无缥缈的传说。

叶牧谦这次没有再用言语刺激，而是安慰道：“天骨之说，缥缈难寻，或许需更大机缘，或许本就是一种理想境界。不必执着，脚踏实地，将魂骨之境锤炼至圆满，根基扎实，未来道途自然宽广。但……”

叶牧谦发觉白言冰面色有些不对：“为师观你眉间有郁气，可是担心家中？”

白言冰沉默良久，终于开口：“师尊明鉴。弟子……弟子实乃皇室出身。近日朝中动荡，弟子忧心如焚，却不敢擅离，怕连累师尊。”

他苦笑：“那日以死相逼拜师，一半是慕道心切，一半……也是想寻个由头，逃离那龙潭虎穴。”

叶牧谦恍然大悟：难怪他对长生执念如此之深，原来是皇室成员。皇室斗争残酷，若无实力，确实朝不保夕。

绝念峰上的小院，终于不再是叶牧谦急于逃离的牢笼，也不再是充满猜忌与算计的戏台。

他不知道这条胡编的路最终通向哪里，也不知道那奇异的“道行反哺”究竟是何原理。

但此刻，他心中已定。

\vspace{2em}

绝念峰上的日子，在叶牧谦真正定下为师之心后，仿佛按下了快进键，却又流淌出一种前所未有的踏实韵律。

白言冰稳固魂骨境界后，并未因无法触及“天骨”而气馁，反而愈发沉静。

每日除了雷打不动的静坐温养骨髓、运转那日渐圆融沛然的内劲外，便是打磨剑术。黑风岭一战的生死感悟被他反复咀嚼，那些原本略显生涩的内劲运用，渐渐化为一种近乎本能的反应。

他练剑时多了几分凝重与精准，偶尔剑锋划过空气，会带起一丝极淡的、温润的破空颤音，那是内劲高度凝聚的表现。

白千羽则正式以弟子之礼侍奉叶牧谦左右。她本就聪慧，又有扎实的医药底子，在叶牧谦那些夹杂着现代常识与玄学词汇的“点拨”下，对气血、经络、人体与自然关系的理解突飞猛进。

她甚至开始尝试将一些修炼中的静心、调息理念，融入对普通病症的调理思路中，虽未直接修炼“问道途”，但精神日益饱满，眸中神采湛然，隐隐有脱俗之姿。

叶牧谦自己，在持续接收着来自白言冰、白千羽境界稳固带来的细微反哺同时，也开始真正思考“问道途”的下一步。

他意识到，自己不能永远停留在根骨和内劲这种相对粗糙的概念上。

既然这条路走得通，且与自己息息相关，那么构建更系统、更深入的体系，不仅是为了教徒弟，或许也是探索自身奥秘的关键。

他结合前世看过的无数修真设定，开始胡诌……不，是推演“问道途”的后续境界。

在白言冰“魂骨”圆满，内劲充盈之后，下一步，或许该是“养气”？

将内劲进一步提炼、转化，与呼吸、与某种更精微的能量——暂时称之为“元气”——结合，温养五脏六腑，开辟下丹田？

他一边在沙地上写写画画，一边随口向白言冰灌输这些更加玄乎的概念，美其名曰“预先感知方向，夯实当前根基”。

白言冰听得如痴如醉，只觉得道途浩瀚，师尊学究天人。白千羽也时常旁听，虽不修炼，却也能从中获得不少启发，用以印证医理。

这一日，午后阳光正好。白言冰在院中空地上缓缓打着一套叶牧谦胡编的、糅合了太极拳架子与五禽戏形态的“锻体拳”，动作松柔缓慢，内劲却如暗流般在体内周流不息。

白千羽则在屋檐下分拣着她前几日下山采买的药材，不时抬头看看兄长，嘴角含笑。叶牧谦坐在石桌旁，面前摊着几张粗糙的树皮纸，上面画满了只有他自己能看懂的符号和关系图，眉头微锁，似乎在苦思某个关窍。

山间岁月静好，仿佛可以一直这样持续下去。

然而，一阵突兀而急促的“扑棱棱”声，打破了这片宁静。一只通体漆黑、唯独喙部呈暗金色的迅鹰，如同箭矢般穿破云层，精准地掠过小院上空，盘旋两圈后，竟径直朝着白言冰俯冲下来！

白言冰拳势一顿，倏然收功，目光锐利地看向那迅鹰。

这鹰他认得，是皇室专门培育用以传递最紧急情报的，非十万火急之事不会动用！而且，这鹰显然受过特殊训练，认得他的气息！

迅鹰稳稳落在白言冰伸出的手臂上，鹰爪上绑着一根细小的赤金色铜管。白言冰面色瞬间凝重，解下铜管，拧开，抽出一卷薄如蝉翼、却质地奇特的绢纸。展开只看了一眼，他的瞳孔骤然收缩，脸色“唰”地一下变得苍白，连手臂都微微颤抖起来。

“哥？”白千羽察觉不对，放下药材快步走来。

叶牧谦也抬起头，心中升起不祥的预感。

白言冰猛地抬头，看向白千羽，又看向叶牧谦，声音干涩而急促，带着一丝不易察觉的慌乱：“父皇……病重！危在旦夕！急诏我们……速速回京！”

他手中的绢纸飘落，上面仅有寥寥数字，却盖着鲜红欲滴的皇帝私玺——“朕疾笃，二皇子、小公主见诏即刻返京，迟恐不及！白隐手书。”

父皇？二皇子？小公主？

叶牧谦脑子“嗡”的一声，虽然早有隐约猜测这对兄妹出身不凡，但“皇室旁支”和“皇子公主”之间的差距，未免也太大了！他看向白千羽，只见她也愣在原地，脸上血色尽褪，嘴唇翕动，却发不出声音，显然这突如其来的身份揭露和噩耗，让她也懵了。

“你们……”叶牧谦张了张嘴。

白言冰深吸一口气，强行压下眼中的震动与焦急，对着叶牧谦快速解释道：“师尊，事出突然，恕弟子此前隐瞒！弟子白言冰，确是大周国二皇子。这是舍妹白千羽，父皇亲封的安平公主。”

此刻，他脸上再无平日里的单纯与执拗，取而代之的是一种属于皇室子弟的凝重与决断，虽然仍显青涩，却已初具棱角。

“我离宫时，曾留信给心腹侍卫，约定若三月不归，便放玄鹰寻我。如今提前触发，父皇病重，诏书急催，我与千羽必须立刻动身！”

白言冰语速极快，“师尊，京城局势复杂，御医束手，父皇之疾恐怕非同寻常。弟子恳请师尊随我们一同入京！师尊见识广博，或……或有一线生机！”

他眼中满是恳求，甚至带着一丝绝望中的希冀。在他心中，师尊是莫测高深的世外仙人，若这世上还有人能救父皇，非师尊莫属！

叶牧谦头皮发麻。

进宫？给皇帝看病？我哪会啊！我只是个写网文的！

可看着白言冰那几乎要崩溃的眼神，想到这少年对自己的全心信任，再想到自己那莫名其妙的反哺或许与这徒弟的命运紧密相连……他咬了咬牙。

“走！”叶牧谦霍然起身，只吐出一个字。此时此刻，退缩已无可能。

事态紧急，三人几乎没有任何时间收拾。白千羽强忍惊惶，将一些可能用到的急救药材胡乱包起。白言冰则吹了一声尖锐的口哨，那玄鹰振翅高飞，显然是传递某种信号。

不到半个时辰，山脚下竟传来了急促的马蹄声，一队约十人、身着便装却难掩精悍之气的骑士奔至，为首之人下马便拜：“卑职禁军校尉，奉密令接应殿下、公主回京！快马已备好！”

显然，这急诏背后，有着周密的安排。

叶牧谦被白言冰不由分说拉上一匹雄健的黑马。他前世只在景区骑过慢悠悠的马，此刻情急之下，也只能硬着头皮抱住马颈。白言冰翻身上马，坐在他身后，低喝一声：“师尊，得罪，抱稳！”

说罢，一夹马腹，骏马如离弦之箭般冲出。白千羽也被一名女侍卫护着上马，紧随其后。十名护卫前后簇拥，马蹄踏破山间寂静，朝着东北方向，风驰电掣而去。

一路换马不换人，除了必要的进食饮水，几乎毫不停歇。叶牧谦被颠得七荤八素，骨头都快散架，但体内那因反哺而增强的体质此刻显出了作用，竟也硬生生扛了下来。

他心中亦是惊涛骇浪：皇子？公主？自己这便宜徒弟的来头，也太吓人了！这趟皇宫之行，怕是龙潭虎穴。

三日后，风尘仆仆的一行人，终于望见了大周皇城的轮廓。那是一座巍峨巨城，城墙高耸，殿宇连绵，在夕阳下泛着恢弘而沉重的光泽。但此刻，无人有心情欣赏。

通过特殊渠道验明身份后，队伍直接从侧门疾驰入城，穿过一道道肃穆的宫门，最终在一片气氛凝重、侍卫林立的重重殿宇前停下。

“二皇子殿下、安平公主到——！”内侍尖细的嗓音在暮色中响起，带着惶急。

白言冰跃下马，也扶下几乎虚脱的叶牧谦，对迎上来的一位老太监和一位与白言冰容貌颇为相似的年轻人急声道：“福公公，皇兄，父皇如何？”

福公公面色悲戚，眼窝深陷，看了一眼奇装异服的叶牧谦，眼中闪过一丝诧异，但不敢多问，连忙躬身引路：“陛下在养心殿，已昏迷数次，御医……御医们都说，让准备后事了……殿下、公主快随老奴来！”

而那年轻人便是当朝太子白承烨，白言冰长兄，喜怒不形于色，颇有些未来帝王的气质。他眉头一皱，继而嘴角微微下沉——那是一个兄长看到顽劣弟弟安全归来时，既松了口气又强压怒意的表情。

他开口就是数落：“这两个月，你哪去了？父皇病重，朝局动荡，你可知有多少人上疏，说你顽劣，甚至有人揣测你勾结外藩？”

白言冰向白承烨示意：“出去学了点本事，这位是我师尊。”

白承烨一脸黑线。他瞅了叶牧谦一眼，满眼都是怀疑。

白言冰的性格，他可是非常知道的：倔强，认准一件事之后，把南墙撞碎了也不回头。白言冰自幼体弱多病，却偏偏抱着一些问道长生的幻想……或许这就是越缺什么、越需要什么吧？

他显然是见过许多方士试图与白言冰打交道，而这些方士实际上都是胡说八道、骗钱来的，最终连调理他二弟的身体都做不到。这些方士最终大多还是因欺君罔上之罪被处理掉了。

白承烨又转向白言冰，问：“你这师尊，他……真没问题吗？”

白言冰顾不上反驳。“这次肯定没问题！”于是冲进养心殿，白千羽紧随其后。

白承烨脚步没有那么快。他再次把目光转向叶牧谦，似乎想在他身上找到常规方士面对此等大场面时应有的慌乱。但叶牧谦此刻却只有好奇，东看看西看看，白承烨一时还找不到他的毛病。

于是白承烨换上了一副储君审视臣下的冷硬面孔，再度开口，这次是对叶牧谦说的：“欺君罔上是要杀头的。”

叶牧谦回过神来，微笑着点了点头。

虽然他目前应该不算是个“修士”，但这几天偷摸内视时发现自己竟然也进入魂骨期了，再怎么说身体条件也和魂骨期的白言冰差不多。至少他还是比较自信，这些凡人应该是杀不掉他的。

白承烨内心也是狐疑起来：气息沉稳如此，难道说真是有些手段在身？

两人前后脚也踏入了养心殿。

养心殿内，药气浓郁，混杂着一种沉闷的、属于病榻的衰败气息。龙床上，一个面容憔悴、两颊凹陷的中年男子闭目躺着，呼吸急促而浅薄，脸色泛着不正常的潮红，嘴唇却是青紫。旁边，几名穿着御医官服的老者愁眉苦脸，低声商议，却都是一筹莫展的颓然。

“父皇！”白言冰和白千羽扑到榻前，声音哽咽。

叶牧谦被白言冰以“随行医者”的名义勉强带进内殿，他站在稍远处，仔细观察着龙床上那位大周皇帝白隐。高热、急促呼吸、肺音粗重、面色潮红、口唇青紫……这症状，怎么越看越像是……

“陛下所患何疾？”白千羽已迅速收敛悲戚，展现出医者的冷静，转向御医问道，声音虽轻，却带着不容置疑的威严。

为首一位白发御医颤声道：“回公主，陛下起初只是风寒咳嗽，后转为高热不退，咳喘剧烈，痰中带血。臣等用尽清热宣肺、化痰平喘之方，如麻杏石甘汤、千金苇茎汤加减，初时稍缓，旋即反复，且愈发沉重。如今邪热壅肺，灼伤肺络，肺气欲绝，已是……已是肺痈重症，药石罔效啊！”

肺痈？叶牧谦心中一动。

古代所说的肺痈，范围很广，但看这症状描述和高热不退的架势，极有可能是感染引起的严重肺炎，甚至可能伴有肺脓肿！

这在没有抗生素的时代，确实是致命的绝症。

白千羽实际上早已跪在床榻边上，不顾礼仪，轻轻搭上白隐的手腕，又小心查看了他的舌苔、眼睑，倾听他的呼吸声。

片刻后，她眉头紧锁，沉声道：“邪热炽盛，盘踞肺腑，已成燎原之势。此时再用寻常清热轻剂，如同杯水车薪。必须先以猛药破开热毒壅滞，给肺气一线生机，再图缓治固本！为何还用这些平和的方子？”

她指向御医案几上刚刚拟好的、一看就是温补为主的药方。

御医们面露难色，一位稍年轻的忍不住道：“公主明鉴，陛下万金之躯，如今正气已虚，若用虎狼猛药，恐……恐陛下龙体承受不住，万一……我等万死难赎其罪啊！”

这话说出了所有御医的心声，不是他们不知道要下猛药，而是不敢！

治好了未必有大功，治死了必定是抄家灭族的大罪！

皇家御医问题就在这里了：皇室成员一旦生了大病，就只能用温和的药物拖着，听天由命。

白千羽气得脸色发白，却也知道这是宫廷痼疾。她下意识地看向叶牧谦，眼中带着求助与询问：“先生……您看？”

刹那间，殿内所有人的目光都聚到了这个刚刚踏入大殿、穿着怪异、一直沉默旁观的年轻人身上。御医们目光惊疑，福公公眼神探究，白言冰则是满含期待。

叶牧谦压力山大。他懂个屁的医术！但此刻箭在弦上，他想起以前自己重感冒发烧，医生开的药里似乎总有鱼腥草成分，说是抗病毒消炎……虽然细菌和病毒不同，但鱼腥草好像确实有清热解毒、消痈排脓的功效，中医常用治肺痈。而且这药……药性应该比较猛烈，味道也冲。

死马当活马医吧！总比等死强！

他心一横，努力回忆着看过的零星中医知识，尽量让自己的语气显得高深莫测：“邪热炽盛，胶结成痈，非峻烈之品不能破其坚垒。徒……白姑娘所言极是，当务之急，是破痈排毒。”

他顿了顿，在众人聚焦的目光中，缓缓吐出三个字：“鱼腥草。”

“鱼腥草？”御医们一愣，随即哗然。

“此物虽可清热解毒，但性寒峻烈，气味腥臭异常！陛下如今脾胃虚弱，如何受得住？”

“正是！且所需剂量必定巨大，这腥臭之气，恐陛下未服先呕，徒伤胃气！”

“荒谬！从未听闻以此为主药治帝王重症者！”

白千羽却眼睛一亮。鱼腥草！她自然知道这味药，其消痈排脓之力确实很强，但正如御医所言，其性猛烈，味道难以忍受，通常只作佐使，或用于民间疾苦、体质强健者。

但此刻父皇肺痈已成，热毒壅盛到御医不敢下药的地步，不正需要这样一味“悍将”来撕开缺口吗？这与她所想“先用猛药祛病救命”的思路不谋而合！

“先生高见！”白千羽脱口而出，随即转向御医，斩钉截铁道：“就用鱼腥草！取其鲜品，绞汁，大量灌服！佐以少量生姜汁护胃，蜂蜜调味。如今已到生死关头，循规蹈矩只有死路一条！一切后果，由我承担！”

“千羽！”白言冰看向妹妹，眼中既有担忧，也有支持。而白承烨的眉头皱得更深了，不知道他在想些什么。

“安平公主，此非儿戏！”老御医劝谏道。

就在这时，龙床上传来一阵微弱却清晰的咳嗽声，白隐竟然悠悠转醒，眼神涣散，却似乎听到了他们的争论，气若游丝地道：“吵……吵什么……朕……朕还没死……千羽……你既敢说……就按你说的……试……”说完，又剧烈咳嗽起来，痰中带着触目惊心的血丝。

皇帝发话，御医们纵然万般不愿，也只能遵旨。只是看向白千羽和叶牧谦的眼神，充满了不忿与等着看笑话的冷意。

新鲜的鱼腥草被紧急寻来，捣烂绞汁。那浓烈刺鼻的腥臭气息瞬间弥漫整个养心殿，不少宫人都忍不住掩鼻。药汁呈暗绿色，看着就让人望而生畏。白千羽亲自接过，混合入少许姜汁与蜂蜜，跪在榻前，柔声唤道：“父皇，服药了。”

白隐勉强睁开眼，闻到那味道，眉头紧皱，喉头滚动，显然极度不适。但他看了一眼眼神坚定的女儿，又看了一眼旁边虽然紧张却同样坚定的儿子，还有那个站在阴影里、看不清表情的“世外高人”，终于闭上眼，任由白千羽将药汁一点点喂入。

药汁入腹不久，白隐果然反应剧烈，开始呕吐，将方才服下的药汁和少许胃内容物都吐了出来，脸色更加难看。接着又开始腹泻。养心殿内一片忙乱，御医们连连摇头，眼神仿佛在说“看吧，闯祸了”。

白千羽脸色苍白，却咬着牙，待白隐稍微平复，又端来第二碗：“父皇，忍住，必须服下去！”

叶牧谦在一旁看得心惊肉跳，手心里全是汗。他只知道鱼腥草可能有用，但具体会怎样，会不会加速死亡，他完全没底。

白言冰紧紧握着拳头，指甲掐进肉里。

白隐在痛苦的呕吐腹泻中，又被灌下了第二碗、第三碗……直到精疲力竭，昏睡过去。

接下来的两日，是无比煎熬的等待。白隐时而昏睡，时而因呕吐腹泻醒来，整个人迅速消瘦下去，气息奄奄。御医们几乎已经认定回天乏术，福公公、白承烨等人即使不懂医术也觉得真没救了，私下里已经开始准备后事。

宫中的气氛压抑到了极点。

白千羽衣不解带地守在榻前，不时为父皇擦拭，调整后续的温和调理方剂，眼神却始终没有绝望。叶牧谦和白言冰也几乎没合眼。

转折发生在第三日清晨。

昏睡中的白隐，那一直紧蹙的眉头，似乎舒展了一些。原本粗重急促的呼吸，变得稍稍平缓悠长。最明显的是，他脸颊上那不正常的潮红，开始减退。

到了午后，白隐竟然自己醒了过来，虽然依旧虚弱，但眼神却清明了不少。他轻轻咳了几声，白千羽连忙递上痰盂，吐出的痰液，虽然依旧浓稠，但血色已淡了许多。

“朕……朕觉得……胸口……没那么堵得慌了……”白隐声音沙哑，却带着一丝久违的轻松。

“父皇！”白千羽喜极而泣。

御医们被紧急召来诊脉，一个个面露不可思议之色。脉象虽然依旧虚弱，但那股沉伏紧促、灼热壅塞的邪气，竟然真的被撼动、被削弱了！肺部的啰音也明显减轻！

“奇迹……真是奇迹！”老御医喃喃道，看向白千羽和叶牧谦的眼神，充满了震惊与复杂。

白千羽抓住时机，立刻调整方药，以沙参麦冬汤、生脉散等方加减，益气养阴，清解余热，固护肺脾。

一日好过一日。高热彻底退了，咳嗽渐止，痰液变清变稀。白隐的胃口也慢慢恢复，虽然只能喝些米汤，但精神明显好转。

不出半月，这位原本被御医判了“死刑”的大周皇帝白隐，竟然已经能被人搀扶着在殿内慢慢行走，脸上也有了血色。虽然离完全康复尚需时日，但性命已然无忧，恢复康泰只是时间问题。

唯一不太高兴的只有白承烨，毕竟如果白隐真去世了，他就能够登临大位了；但毕竟为人子女，也为父亲的康复感到由衷的喜悦。

这一日，天气晴好。白隐坐在铺着厚垫的椅子上，晒着太阳，看着侍立在一旁、气质已然大变的儿子，和正在仔细为他调整药膳配比的女儿，心中感慨万千。

他看向白言冰，本想板起脸教训几句：“言冰，你说是出去登山游览，散散心，结果一去两月不回，音讯全无。你这些日子，到底跑哪里去胡闹了？知不知道朕和你母后有多担心？”

但话到嘴边，却咽了回去。因为他清晰地感觉到，眼前这个儿子，和两个月前离开时，已经判若两人。不仅仅是身形更加挺拔、气质更加健康；更有一种由内而外散发出的沉凝、自信、仿佛经历过生死淬炼的气质。

那眼神清澈坚定，却又深邃宁静，再无往日提到“仙道”“长生”时的那种虚浮狂热，取而代之的是一种脚踏实地的、内敛的力量感。

这变化太大，大到让白隐觉得，任何关于“胡闹”的斥责，都显得苍白无力。

最终，白隐只是深深看了白言冰一眼，化作一声复杂的叹息，摆了摆手：“罢了……回来就好。看来，你这趟出去，并非一无所获。”

他的目光又转向不远处似乎与宫廷氛围格格不入的叶牧谦。此时白承烨站在他面前，因先前的错误怀疑，正拱手向他道歉呢；叶牧谦则扶起他，说着什么“不知者不为罪”之类的话。

这个年轻人，穿着依旧古怪，但经过此事，白隐再不敢以常理度之。正是他轻描淡写的一句话，点出了“鱼腥草”，才让自己从鬼门关捡回一条命。而女儿千羽对他的态度，更是恭敬有加，甚至超越了对父皇的依赖。

“叶……先生。”白隐斟酌着称呼，“朕这条命，多亏了你和千羽。救命之恩，朕记下了。”

叶牧谦连忙躬身：“陛下洪福齐天，草民不敢居功。是安平公主医术精湛，决断果敢。”

白隐不置可否，笑了笑，又看向白千羽，温言道：“千羽，这次多亏了你。你的医术见识，已远超宫中御医。看来，你离宫这些时日，也另有奇遇。”

白千羽放下手中的药匙，走到白隐面前，郑重跪下：“父皇，女儿能有所进益，全赖先生点拨教导。先生虽未正式收女儿为徒，但传授医理、指点迷津，实同师授。

“如今父皇康健，女儿心愿已了。恳请父皇允准，女儿愿正式拜入叶先生门下，追随先生修习大道！”说罢，她转向叶牧谦，以额触地，行了一个标准的拜师大礼。

殿内一片安静。白隐看着女儿坚定的眼神，又看看儿子眼中对叶牧谦毫无保留的信赖，最后目光落在叶牧谦身上。这个神秘的年轻人，救了自己的命，更似乎彻底改变了自己的一双儿女。

他沉吟片刻，缓缓点头：“叶先生乃世外高人，你能拜入先生门下，是你的造化，朕……准了。”

“谢父皇！”白千羽声音带着欣喜，再次向叶牧谦叩首：“弟子白千羽，拜见师尊！请师尊收录！”

叶牧谦看着跪在面前的少女，心中百感交集。从最初的怀疑监视，到后来的半信半疑，再到如今的全心信赖与正式拜师……这一路，可谓曲折。他伸手虚扶：“起来吧。既入我门，当勤修不辍，坚守本心。”

“弟子谨遵师命！”白千羽起身，眼中光华流转，那是找到了前路的明澈与坚定。

白言冰也为妹妹感到高兴，嘴角露出笑容。

\vspace{2em}

没高兴几日。

这份短暂的宁和，很快被边关一道加急军情撕裂。

“报——！晋国边军异动，于北境苍云关外增兵五万，游骑频繁越境挑衅，边军告急！”

军情奏报如同一盆冰水，浇在刚刚缓过气来的大周朝廷头上。养心殿内，刚刚能起身视事的白隐，看着手中紧急军报，眉头紧锁；原本因病而显瘦削不少的脸上，此刻却浮动着罕见的凝重与锐利。

他生性宽仁，甚至被一些近臣私下认为软弱，但这位帝王并不糊涂。

“朕重病垂危，若非千羽，此刻恐怕已宾天了。此事在宫中亦属绝密，晋国如何知晓？增兵时机，拿捏得如此‘恰好’？”

白隐手指轻叩御案，声音低沉，目光扫过殿中几位重臣，最后落在垂首恭立的宰相赵文渊身上，停留一瞬，又淡淡移开。

“边境不可不防。传旨，调拨京城禁军左卫、右卫，即刻开赴苍云关，增援边军。京城防务……尚有三千御林军，足以维持治安，震慑宵小。”

“父皇！”白言冰踏前一步，躬身请命，“儿臣愿随军出征！定不让晋国越境半步！”

白隐看着气质迥异的儿子，心中欣慰，却缓缓摇头：“你刚回京不久，且留在朕身边。京城……也需要人坐镇。”他语气温和，却带着不容置疑的决断。

一方面，他确实想将这不凡的儿子留在近前；另一方面，潜意识里，他总觉得这突如其来的边患背后，似乎缭绕着更深的迷雾。让白言冰留在京城，或许更能应对变局。

白言冰还想再言，但见父皇神色坚决，只得按下话头，躬身应道：“儿臣遵命。”

数日之后，大军开拔，铁甲铿锵，带走了京城大半精锐，也带走了无数人心中那份踏实。皇城仿佛一下子空旷冷清了许多，只剩下三千御林军鲜明的身影，在宫墙内外规律巡弋，平添几分肃杀。

白千羽已正式随叶牧谦学习，暂居宫内一处清静偏殿。叶牧谦也被白言冰以“客卿”名义安置在靠近皇宫的一处馆舍，待遇优渥，却也处于某种半监视之下。

叶牧谦本人对朝堂风云毫无兴趣，每日除了接受白言冰、白千羽的请教，便是继续闭门推演他那套越来越庞大的问道途体系，同时默默感受着自身随白言冰境界稳固而持续缓慢增强的身体素质。

京城繁华，于他而言，远不如绝念峰清净，但既已卷入，也只能静观其变。

他并未察觉，一张阴谋的大网，已在京城最深处悄然织就。

宰相赵文渊，三朝元老，门生故吏遍布朝野，表面恭顺勤勉，内心却早已被权欲与野望腐蚀。皇帝病重垂危的消息，正是他通过宫中眼线秘传晋国，意图里应外合，搅乱局势。

如今白隐虽然奇迹般康复，但京城空虚，在他看来，仍是千载难逢的时机！

“禁军已远赴边关，京城只剩三千御林军。白隐病体初愈，太子一无用书生罢了，二皇子不过是个痴迷怪力乱神的黄口小儿，安平公主……天赐良机！”

赵文渊在密室中，对几名心腹武将和重金聘来的江湖魁首低语，眼中跳动着贪婪与狠厉的火光。“今夜子时，以火起为号，尔等率家丁私兵，并诸位豪杰麾下好手，直取皇城！先登者，千金赏，万户侯！”

他筹谋已久，暗中纠集了能掌控的武官部曲、蓄养的死士家丁，加上用重金和许诺招揽的所谓“武林高手”，竟也凑出了三千余亡命之徒。

其中不乏在江湖上名头响亮、心狠手辣之辈，如掌法凌厉的“血手人屠”厉刚、善用暗器的“追魂扇”莫先生等，皆被赵文渊许以高官厚禄，网罗麾下。

月黑风高，杀机暗藏。

子时将至，皇城东北角一处偏僻仓房果然冒出浓烟火光，旋即喊杀声四起！无数黑影从各个街巷涌出，手持利刃，呼喝着冲向皇城各门！

御林军虽有所戒备，但叛军来得突然，人数众多，且其中混杂着真正的江湖好手，武艺高强，手段狠辣，绝非寻常乌合之众可比。

“火起！敌袭——！死守！”御林军统领带着惊怒的喊叫在夜空中久久回荡。

刀剑碰撞声、惨叫声、喊杀声瞬间打破了皇城的寂静。叛军有备而来，携带了简易的攻城梯和撞木，加之部分守军内应暗中打开缺口，战斗从一开始就异常惨烈。

御林军虽训练有素，但猝不及防下，阵线被多处突破，叛军如同潮水般涌入第一道宫墙，向内廷宫城逼近。火光映照着厮杀的人影，鲜血染红了汉白玉的台阶。

福公公慌忙闯入养心殿：“启禀皇上，宰相造反啦！”

白隐被惊醒，连衣服都顾不得换，就在福公公和少量侍卫的护持下通过密道快速转移；白承烨也在随从的护卫下撤离。

殿外传来的喊杀声越来越近，御林军节节败退的消息不断传来。所幸内廷宫城尚未陷落，兵士只能死守宫城，但在现在的情况下，宫城陷落也只是时间问题罢了。

“赵文渊！好贼子！”白隐一边沿密道逃跑，一边气得浑身发抖，更多则是心寒。

他自问待赵文渊不薄，却换来如此狠毒的背叛。

“快去叫人，向直隶省厢军求援！”

“父皇！”白言冰疾步跑来，他已迅速披挂上一件轻甲，手中握着的正是那柄从绝念峰带回、饮过黑风岭贼血的佩剑。

他脸上只有一种冰冷的沉静，眼神锐利如即将出鞘的剑锋。

“叛军已攻破皇城神武门，正向内廷杀来。御林军正在浴血奋战，但叛军中有不少倒戈武官、武林高手，抵挡艰难。”

“怎么办……言冰，你可有办法？”白隐此刻只能将希望寄托在这个脱胎换骨的儿子身上。

白言冰深吸一口气，望向殿外冲天的火光和越来越清晰的厮杀声，缓缓道：“擒贼先擒王，挫其锐气。儿臣请命，出战。”

“不可！”白隐下意识反对，“叛军势大，你一人如何能挡？敌众我寡，朕已命人紧急向直隶省求援，即使短时间内放弃皇城、转移去直隶……”

“父皇，援军到来需要时间。此刻叛军气势正盛，若不能当头棒喝，御林军士气恐彻底崩溃，内廷不保；若是轻易放弃内廷，还不知道那赵老头会做出什么龌龊事来！”

白言冰语气坚定，“师尊传授之道，正在实战中砥砺。请父皇准儿臣一试！”

看着儿子眼中那不容置疑的自信与决绝，想起他提及黑风岭时那轻描淡写的“全歼”，白隐心中挣扎，最终咬牙道：“……小心！若事不可为，立刻退回！”

“儿臣领命！”

白言冰转身，大步踏出养心殿密道。

殿外，硝烟弥漫，火光摇曳，远处兵刃交击与惨呼之声不绝于耳。

白言冰提剑，沿着汉白玉的御道，迎着厮杀最激烈的方向，一步步走去。步伐沉稳，不快，却带着一种奇特的韵律，仿佛每一步都踏在心跳的节拍上。

周身那温润如玉的魂骨隐隐散发出微光，体内沛然流转的内劲如同蛰伏的江河，随时准备奔涌咆哮。

前方，正是通往内廷的最后一道重要门户。此处御林军残部正依托门楼和宫墙，与潮水般涌来的叛军殊死搏杀。

叛军为首一员大将，身材魁梧，满脸虬髯，手持一柄沉重的开山斧，正是赵文渊笼络的一名边军叛将，有“开山虎”之称的猛将胡贲。他此刻杀得兴起，一斧将一名御林军校尉连人带盾劈成两半，鲜血溅了他满脸，更显狰狞。

“哈哈哈！御林军不过如此！儿郎们，杀进去，活捉白隐者，赏千金！”胡贲狂笑，声震夜空。

就在此时，他忽然感到一股异样。原本激烈厮杀的前方，自动分开了一条通道。

一个身穿轻甲、手持长剑的年轻身影，独自一人，缓缓走了出来。火光映照下，那人面容俊朗，眼神平静得可怕，仿佛眼前不过是自家庭院。

“哪来的毛头小子？找死吗？”胡贲一愣，随即嗤笑，他根本没把白言冰放在眼里。

观其步伐，不像有什么高深武功；年纪又轻，估计是哪个不知死活的皇室子弟出来充英雄。

白言冰在胡贲身前十步处站定，目光扫过眼前黑压压的叛军和凶神恶煞的胡贲，声音平淡，却清晰传入每个人耳中：“逆贼犯阙，罪不容诛。现在放下兵器，还可留个全尸。”

“哈哈哈哈！”胡贲和周围叛军爆发出更大的哄笑，仿佛听到了天大的笑话。

“小娃娃，乳臭未干，也学人阵前喊话？”胡贲狞笑着，将开山斧往地上一顿，砸得石板碎裂。

“爷爷我杀人时，你还在吃奶呢！来来来，爷爷陪你玩玩，一对一打，免得说我欺负你！”他存了戏耍和震慑的心思，打算亲手虐杀这个不知天高地厚的小子，进一步打击御林军士气。

白言冰没有废话，只是微微抬起了手中的剑，剑尖斜指地面。

胡贲大吼一声，如同猛虎出闸，沉重的开山斧带着凄厉的风声，一招最简单粗暴的“力劈华山”，朝着白言冰当头劈下！这一斧凝聚了他全身蛮力，足以开碑裂石，寻常武者别说硬接，光是那气势就足以夺人心魄。

然而，在白言冰的感知中，这一斧虽猛，却轨迹清晰，力量涣散，破绽百出。他脚步未动，只是在那斧刃即将临头的刹那，身体以毫厘之差微微一侧。

巨斧带着狂风，擦着他的甲叶劈空，重重砍在地上，碎石飞溅。

胡贲一愣，力道用老，正要回斧变招，却见眼前剑光一闪！

那剑光并不璀璨，甚至有些朴素，只是快！快得超越了胡贲神经反应的速度！快得仿佛它本来就在那里！

一道温润却凌厉无比的剑气悄然划过胡贲粗壮的脖颈。

胡贲脸上的狞笑瞬间凝固，眼中充满了难以置信的惊愕。他张了张嘴，想说什么，却只发出“嗬嗬”的漏气声。

下一刻，斗大的头颅冲天而起，鲜血如喷泉般从无头的脖颈腔子里狂涌而出，魁梧的身躯晃了晃，便轰然倒地。

一回合！仅仅一个照面，叛军骁将“开山虎”胡贲，授首！

乾元门前，出现了刹那的死寂。无论是叛军还是御林军，都被这电光石火、干净利落到极致的一剑惊呆了。那看似轻描淡写的一侧身、一递剑，却蕴含着令人胆寒的精准、速度与力量！

白言冰甩了甩剑尖上并不存在的血珠，目光平静地投向呆若木鸡的叛军。

他一个字都没有说，身影却已然消失在了原地。

下一瞬，他如同鬼魅般撞入了叛军最密集之处！

这一次，不再是单挑，而是真正的大开杀戒！

他没有使用任何花哨的剑招，只是最简单的刺、削、抹、点。但每一剑都妙到毫巅，快如闪电。剑锋所向，叛军手中的刀枪如同纸糊般被轻易挑飞、斩断。

附着内劲的剑身，无需刻意劈砍；这些仅是凡人的叛军，挨着些就死，擦了边就亡。

白言冰步法飘忽，在人群中穿梭自如，叛军密密麻麻的刀枪往往连他的衣角都碰不到，而他的剑却总能从最不可思议的角度，带走一条条性命。

更可怕的是他的防御。偶尔有冷箭或偷袭的兵器及身，却只能在那看似普通的轻甲甚至露出的肌肤上，留下一道浅浅白痕，发出金铁交鸣之声，根本无法破防！

魂骨之境，淬炼周身骨骼如玉如钢，衍生出的内劲自动护体，岂是寻常凡人弓弩所能伤的？

他就像一台精密而高效的杀戮机器，所过之处，人仰马翻，血花绽放。叛军起初还仗着人多试图围攻，但很快就变成了单方面的屠杀。惨叫声、骨骼碎裂声、兵器折断声不绝于耳，中间夹杂着白言冰那稳定到令人心寒的脚步声和剑锋破空的轻微颤音。

“怪……怪物啊！”

“他不是人！快跑！”

叛军的勇气和阵型，在这非人的战力面前迅速崩溃。白言冰一个人，竟生生在叛军潮水中，犁出了一条血肉胡同！他所向之处，叛军无不胆裂，纷纷溃散。

原本士气低落、苦苦支撑的御林军，目睹这神迹般的一幕，先是目瞪口呆，随即热血上涌，狂喜与震撼交织！

“二殿下神威！”

“跟着殿下，杀逆贼！”

绝境之中，白言冰以一人一剑，力挽狂澜，不仅瞬间斩杀了敌将，更以碾压般的实力横扫叛军前锋，极大地迟滞了叛军的攻势，为御林军赢得了宝贵的喘息时间。御林军统领趁机收拢残部，重整旗鼓，原本濒临崩溃的防线迅速稳固下来。

“放箭！掩护二殿下！”统领嘶声下令。

箭雨倾泻，将试图重新组织进攻的叛军压了回去。御林军将士们眼含狂热，紧握兵器，跟随着前方那道如同战神般的身影，发起了反击。

士气此消彼长，战局瞬间扭转！

叛军后阵，赵文渊在几名高手护卫下，远远望见胡贲被秒杀、前锋被一人杀得溃不成军的情景，惊得面如土色。“那……那是白言冰？他……他何时有了这等武功？！”他心中骇然，隐隐感到不妙。

“相爷勿忧！不过是些粗浅横练功夫加上速度快些罢了！待我等前去，取他首级！”厉刚舔了舔嘴唇，眼中露出嗜血的光芒。他自恃武功高强，杀人无数，觉得白言冰和他这种人差不多，只是仗着身法和力气欺负普通士卒。

“不错，我等联手，必能杀他！”“莫先生摇着铁扇，阴恻恻地道。

顿时，七八名被重金聘来的、在江湖上颇有名号的“高手”越众而出，各持兵刃，狞笑着向白言冰包围过去。他们身形矫健，步伐灵动，显然与普通叛军不可同日而语。

白言冰刚刚一剑荡开三名叛军的长矛，顺手点倒两人，便觉数道凌厉气息从不同方向锁定了自己。他停下脚步，持剑而立，目光平静地扫过围上来的这些江湖客。

从他们鼓荡的气息和站姿来看，比黑风岭那几个头目和要强上不少，大概相当于他外劲大成、内劲初生的水准？在他魂骨圆满、内劲圆转如意的感知中，这些人的气息强弱，甚至下一步的意图，都隐隐有所映照。

白言冰不禁眯着眼睛看了他们一眼，目光中满是蔑视。

而这种蔑视反而使得那些自视甚高的“武林高手”们越发气恼了。

“小子，杀几个废物就以为天下无敌了？”厉刚双掌赤红，散发着淡淡腥气，显然掌功带毒，“让你见识见识什么叫真正的武功！”

说罢，他率先发难，身形如鬼魅般飘忽逼近，一双血手带起腥风，直拍白言冰胸腹要害，掌风未至，一股令人作呕的甜腥气已扑面而来。

与此同时，莫先生的铁扇张开，扇骨中机簧响动，数点寒星疾射白言冰后心与双腿。侧方，一名使细剑的刺客悄无声息刺向肋下，另一名持鬼头刀的壮汉则大吼一声，势大力沉地拦腰横斩！

几人配合默契，封死了白言冰所有闪避空间，显然常做这等以多欺少的勾当。

远处观战的御林军和叛军都不由屏住了呼吸。这几人任何一个放在江湖上都是好手，联手之下，威力更增！

面对这足以让普通武林高手手忙脚乱的围攻，白言冰眼神依旧古井无波。他甚至没有去看那些袭来的暗器和兵刃，只是在内劲感知中，清晰地“看”到了每一道攻击的轨迹、力度，以及它们之间那细微的、因配合并非天衣无缝而产生的空隙。

他动了。

没有后退，反而向前踏出半步，正好切入厉刚双掌力道将发未发、旧力略滞的那一瞬空隙。手中长剑以一种玄妙的弧度向上斜撩，剑身震颤，发出一声清越的嗡鸣，内劲勃发！

“铛！”

剑尖精准无比地点在厉刚右手腕脉之上，一触即收，一股凝练如针的内劲瞬间透入！厉刚只觉整条右臂一麻，赤红的掌劲瞬间溃散，经脉如遭针刺，惨哼一声，踉跄后退，右掌已软软垂下。

白言冰身形借着这一撩之势微微旋转，仿佛早有预料般，剑锋顺势划过一个半圆。“叮叮叮”几声轻响，莫先生射来的淬毒暗器竟被剑脊精准拍飞。旋转未止，他左脚为轴，右脚如鞭子般向后踢出，正中那使细剑刺客的手腕，“咔嚓”骨裂声清晰可闻，细剑脱手飞出。

此时，那鬼头刀壮汉的横斩才堪堪及身。白言冰刚刚踢出的右脚落地，重心转换如行云流水，持剑的右手手腕一翻，长剑由斜撩变为直刺，后发先至，内劲外显，剑尖吞吐着微不可察的淡金色毫芒，点向壮汉咽喉。

壮汉大惊，想要回刀格挡已然不及，只能拼命侧头躲避。

“噗嗤！”

剑尖虽未刺中咽喉，却点在了壮汉的肩井穴上。同样是一触即收，但内劲透入，壮汉整条左臂连同半边身子瞬间酸麻无力，鬼头刀“哐当”落地，庞大的身躯失去平衡，轰然坐倒。

兔起鹘落，不过呼吸之间。围攻的四名“高手”，一人腕脉被废，一人手腕骨折，一人肩颈被制，丧失战力。

其余几人攻势方才到位，却已失去了目标和配合，反而显得笨拙可笑。

白言冰身影再闪，如同穿花蝴蝶，在剩余几人惊骇的目光中，剑光连闪。

“啊！”“我的手！”“我的腿！”

惨叫连连，兵刃坠地声不绝。这些在江湖上也算一方豪强的“高手”们，在白言冰面前，竟如同孩童般不堪一击！

他们的招式、经验，在一名对身体绝对掌控、对内劲运用圆转、以及对战斗节奏和破绽完全洞察的真真正正的修士面前，显得如此粗糙和缓慢。

这远远不是技术的差距，而是单纯的境界碾压！

当最后一名“高手”被白言冰用剑身拍中后脑，哼都没哼一声软倒在地时，叛军阵营中，最后一点抵抗的勇气也彻底消失了。连重金请来的“高人”们都像砍瓜切菜一样被放倒，他们这些普通人还打什么？

“逃啊！”

“二皇子是神仙下凡！快跑！”

叛军彻底崩溃，丢盔弃甲，四散奔逃。

御林军士气如虹，乘胜追击，配合白言冰的余威，很快将残余叛军分割包围，逐一剿灭或擒拿，作乱的武官也大多被当场格杀或擒获。

宰相赵文渊见大势已去，面如死灰，在亲信死士护卫下试图趁乱逃出京城。却没想到，当他们试图出门时，正撞上那个身穿轻甲、手持长剑的身影，他的后面还跟着百十名御林军。

“老赵，跟我走一趟吧。”白言冰面无表情地说。

赵文渊自知不能抵抗，只得束手就擒。其余亲信死士等见赵文渊都主动就缚，也自杀的自杀、投降的投降了。

于是，一场险些颠覆皇权的叛乱，竟因白言冰一人之力，被迅速平定。从白言冰提剑出战，到叛军彻底崩溃，前后不过半个时辰。

天色微明时，厮杀声已然平息。

皇城内外，到处都是战斗的痕迹和叛军的尸体，御林军正在打扫战场，押解俘虏。空气中弥漫着浓重的血腥味和烟熏火燎的气息。

原来昨夜白言冰去后，白隐因不放心儿子，于是不顾阻拦，与侍卫们又返回了养心殿，亲眼目睹了后半程的战斗。

他看到了儿子如何一剑斩将，如何如入无人之境般在叛军中杀戮，如何轻描淡写地击败那些所谓的“武林高手”……

整个过程，白言冰都显得那么从容，那么举重若轻。

白隐的心中，充满了难以言喻的震撼。

这真的是他那个曾经体弱多病、只知痴迷仙道传说、甚至有些迂腐的儿子吗？

这分明是一尊降临凡间的战神！是真正掌握了超凡力量的修行者！

叶先生……那位神秘的叶先生，究竟传授了何等可怕的法门？！

当白言冰提着滴血长剑，踏着晨曦微光，一步步走回养心殿前时，他身上轻甲沾染了血迹，脸上却依旧平静，只有眼中残留着一丝战斗后的锐利精光。

他对着白隐躬身一礼：“父皇，叛乱已平，逆首赵文渊已被擒获。儿臣复命。”

白隐看着儿子，嘴唇动了动，原本心中残存的、对于儿子“不务正业”去修炼的那点芥蒂和担忧，此刻烟消云散。

实际上直到此刻，白隐才真正认识到，白言冰的眼界与实力，早已飞升到了他无法理解、更无法约束的层次。他所学到的，是真正的“真东西”。

他深吸一口气，上前亲手扶起白言冰，用力拍了拍他的肩膀，声音带着激动和后怕，更带着无比的欣慰与骄傲：

“好！好！好！吾儿真乃神人也！此番若非我儿，这大周江山怕不是早倾覆了，此战注你头功！今日起，你想做什么，便去做！父皇再不干涉你了！”

白言冰感受到父皇发自内心的认可与放手，心中也是一暖，再次行礼：“谢父皇。”

就在白隐为儿子的神勇而大喜，准备下令彻查叛党、论功行赏之际，被押解经过殿前、面如死灰的宰相赵文渊，在瞥见静静站在稍远处、似乎与这一切格格不入的叶牧谦时，眼中闪过了一丝极为怨恨的神色。

\vspace{2em}

平叛后的第三日，面色依旧有些苍白的白隐，亲自来到刑部天牢，提审了已成阶下囚的宰相赵文渊。

烛火摇曳，映照着一夜之间仿佛老了十岁的帝王，和虽身陷囹圄、眼中却仍残留着不甘与怨毒的昔日权臣。

白隐的声音低沉而冰冷，带着压抑的怒火：“赵文渊，朕自问待你不薄，三朝元老，位极人臣。为何行此大逆不道之事？你与晋国，究竟有何勾结？璟妃……陈璟当年病逝，是否也与你有关？”最后一个问题，问得异常艰难。

陈璟，那个温婉如水、他真心喜爱却红颜薄命的女子，是白千羽的生母，也是他心中一处柔软的隐痛；虽然陈璟是“暴病而亡”，但白隐觉得她的死绝对不可能那么简单；赵文渊被擒前瞥向叶牧谦那复杂难言的一眼，让白隐心中的不安又加重了几分。

或许是知道自己绝无活下来的可能，或许是积压多年的怨毒与秘密终于找到了宣泄的出口，赵文渊忽然发出一阵嘶哑而癫狂的低笑。

“待我不薄？哈哈……陛下，臣侍奉三代君王，呕心沥血，到头来也不过是个高级些的奴才！凭什么你白家就能世代坐拥江山，我赵文渊就只能俯首称臣？”

他眼中闪烁着疯狂的光，“不错，晋国边患是我通风报信！陈璟……那个蠢女人，也是我安排的！”

白隐瞳孔骤缩，放在膝上的手猛地握紧。

赵文渊似乎很享受皇帝震惊的表情，继续用那种令人毛骨悚然的语气说道：“早在二十年前，你还是个黄毛小子的时候，我就开始布局。陈璟，本是我精心培养，容貌、才情、性情无一不是上上之选；送入宫中，就是为了固宠，诞下带有我赵家血脉的皇子！

“届时母凭子贵，我再徐徐图之……可惜，天不助我！”

他脸上露出扭曲的恨意，“那贱人入宫前说什么都不让我碰，我只能下药让她昏过去，才终于让她怀有我的骨肉。我赌是男胎，结果……却生下了白千羽这个没用的丫头片子！

“不过女流之辈倒也无所谓了，实在不行我冒天下之大不韪、扶持她当女帝，这江山也是我赵家的！后来陈璟竟然察觉了我的意图，竟想向陛下告发……”

每一个字，都如同淬毒的冰锥，狠狠扎进白隐的心脏。他最宠爱的女儿，竟然不是自己的骨血！他真心爱过的妃子，竟是一场处心积虑的阴谋，甚至她的死……也是因为发现了这肮脏的秘密！

但赵文渊毫无停下的意思，反而更加疯狂：“黑风岭山贼，也是我放出来的！我本来探听到白言冰和白千羽的去向，想借这伙贼人把白言冰和白千羽直接就留在那里，没想到出了叶牧谦那个变数！”

“你……你这逆贼！畜生！”白隐猛地站起，浑身颤抖，指着赵文渊，脸色铁青，胸口剧烈起伏，几乎要吐出血来。

他一生宽仁到有些懦弱，甚少动怒，但此刻的暴怒与心痛，几乎要冲垮他的理智。

“来人！给朕……给朕……”

他想说拖出去凌迟处死，但最终残存的帝王心术让他强行压下了立刻处死的冲动。

此事关乎皇室颜面，关乎璟妃清誉，更关乎千羽那孩子的命运！

他深吸几口气，勉强稳住身形，声音从牙缝里挤出：“此事，在场无论是谁，若有半句泄露，朕诛你九族！”

\vspace{2em}

白隐的本意是好的。他想：若是捂住这个秘密，那么白千羽就永远是他最疼爱的小公主，璟妃在史书和世人眼中也依然是那个温婉贤淑的妃子，皇家颜面也得以保全。

他甚至开始思考，如何在不惊动任何人的情况下，给千羽更多的补偿和关爱。

然而，皇宫从来不是密不透风的墙。

赵文渊经营多年，党羽虽遭清洗，但仍有漏网之余孽。这惊天的丑闻，不知通过哪个环节，如同瘟疫般悄然泄露了出去。或许正是那些不甘失败的宰相余党，意图在最后时刻搅乱朝局，报复皇家。

起初只是宫闱深处最隐秘的流言，但很快便如同长了翅膀，飞越宫墙，在京城某些特定的圈子里悄然传播。

当白隐惊觉时，流言已如野火。

白隐大怒，半是为了灭口，半是为了泄愤，把宰相余党满门抄斩，但在不少人眼中反而更是坐实了这些谣言。

“听说了吗？安平公主……并非陛下亲生！”

“是先璟妃与逆臣赵文渊的……唉，真是……”

“皇家血脉岂容混淆？此乃滔天大罪啊！”

“难怪二皇子与公主容貌不甚相似……”

这些窃窃私语，最终不可避免地，传到了当事人耳中。

这一日，白千羽正在自己宫中一处僻静的暖阁内静坐调息。自正式拜师叶牧谦后，她虽未直接修炼与战斗直接相关的外劲法门，但在叶牧谦的指点下，也开始尝试以医入道，调和自身气血，温养精神，并隐隐感知自身根骨状况，为将来可能的修炼打下基础。

她天资聪颖，进展颇快，近日正觉气机活泼，隐隐有触及魂骨门槛的迹象，便更加勤勉，每日抽出大量时间静修。

窗外忽然传来两名洒扫宫女压得极低、却因距离太近而清晰可闻的议论声：

“……真的假的？公主她……”

“嘘！小声点！不要命了？不过……好像是真的，我有个同乡在刑部当差，听他说赵相死前亲口招供的……”

“那……那公主岂不是……唉，真可怜，陛下对她那么好……”

“好有什么用？不是皇家血脉，这公主之位，甚至性命都怕是悬了。听说有些大人都已经准备上奏了……”

字字句句，如同最锋利的冰锥，瞬间刺穿了白千羽沉浸在宁静修炼中的心湖。

不是……亲生？

赵文渊的女儿？

母妃……是被害死的？

无数破碎而尖锐的信息，伴随着难以置信的震惊、被欺骗的愤怒、对身世飘零的恐惧、对母妃枉死的心痛……种种极端负面情绪，如同决堤的洪水，轰然冲垮了她勉强维持的平静心绪。

“噗——！”

气血逆行，内息瞬间紊乱！她正在尝试感应、引导的体内那点微弱气机，此刻被狂暴的情绪裹挟，失去了控制，如同脱缰野马，在她经络中横冲直撞！

“问道途”的修炼，重意重悟，进展极快，但对心境要求也极高。叶牧谦胡编时未曾细想，但这套基于他潜意识里无数小说设定拼凑的体系，无形中暗合了某种高风险的特性。

心念纯粹、意志坚定如白言冰，尚需生死搏杀来淬炼心志；而白千羽此刻遭遇的，是颠覆自我认知、撕裂情感根基的剧变，其冲击力远超寻常挫折。

心魔，由此滋生！

“啊——！！！”

一声凄厉不似人声的尖叫从暖阁中爆发。紧接着是器物被狠狠砸碎的刺耳声响，桌椅翻倒的碰撞声，以及宫女惊恐的哭喊和逃窜声。

暖阁内，白千羽双目赤红，原本清丽温婉的脸庞扭曲着痛苦与狂乱，长发披散，周身气息狂暴紊乱，竟将靠近的桌椅摆设凭空震飞、撕裂！

她仿佛失去了神智，只剩下本能地发泄着体内无处安放的痛苦与力量，所见之物，尽皆成为摧毁的对象。

“公主！公主您怎么了？快来人啊！”

宫人们吓得魂飞魄散，无人敢上前。

几乎是同一时间，正在宫外馆舍中推演养气境、试图理清“元气”与内劲关系的叶牧谦，心头猛地一跳。

一种极为不安的悸动感传来，并非通过声音，更像是一种冥冥中的感应——与他道途相关的弟子，出现了巨大的危机！是白言冰？不对，白言冰的气息沉凝依旧。是……千羽！那种气息紊乱、近乎走火入魔的征兆！

他霍然起身，来不及细想，便朝皇宫方向疾奔而去。然而，到了坤宫后门，却被一名面色冷硬的御林军军官拦下。

“站住！皇宫重地，无诏不得擅入！叶先生，即便是公主师尊，也需通传获准！”军官手按刀柄，语气不容置疑。

叛乱刚过，宫禁森严到了极点。

“让开！白千羽有危险，我必须立刻进去！”叶牧谦急道。

“规矩就是规矩，你是她老师也不行！再说，你一个男人，进入后宫，是何居心？”军官寸步不让，身后士兵也警惕地围了上来。

原来，为了防止秽乱宫廷，大周国宫城分为两部分。两宫之间只有皇帝、宫女和太监能自由进出，否则除非有谕令，两宫之间是不能人员交流的；只是白隐特别宠爱白千羽，允许她出来玩罢了。

叶牧谦虽显然是有些神秘手段，但此刻显然不能硬闯皇宫，那只会让事情更糟。

他心急如焚，却无可奈何，只能对那军官吼道：“快去禀报二皇子白言冰！立刻！告诉他千羽出事了！”

那军官虽然认死理，但也给旁边士兵使了个眼色。士兵们对视一眼，开始向门内高声喊叫：“快去禀报二皇子殿下！”门内守着的太监宫女听令，立即奔去乾宫。

白言冰其时正在自己宫中清修，同时回味着父皇那句“再不干涉”所带来的轻松感。忽然接到宫中急报，言公主宫中突发异状，似有走火入魔之象，叶先生被阻宫外，情势危急。

白言冰的脸色瞬间变了。千羽走火入魔？他来不及细想原因，身形已如离弦之箭般射出，朝着坤宫狂奔。

什么后宫礼仪，什么男女大防，此刻全被抛诸脑后。他心中只有一个念头：千羽不能有事！

他速度极快，宫门门监只见一道青色残影掠过，甚至来不及阻拦，人已闯入深宫。循着那越来越清晰的混乱与狂暴气息，他径直冲到了白千羽所在的暖阁之外。

只见阁内一片狼藉，木屑瓷片遍地，白千羽状若疯魔，仍在无意识地挥掌踢腿，道道紊乱的气劲将室内破坏得更加彻底，几名瑟瑟发抖的太医和宦官躲在远处，不敢靠近。

“千羽！”白言冰大喝一声，声音中灌注了一丝凝实的内劲，试图唤醒她的神智。

然而，陷入心魔的白千羽恍若未闻，反而被声音刺激，赤红的双眸猛地锁定白言冰，发出一声低吼，竟合身扑上，一掌拍来！

掌风凌乱，却带着一股不顾一切的狠戾。

白言冰心中剧痛，不敢硬接，更不敢伤她，只能施展身法轻盈闪避，同时不断呼唤：“千羽！是我！哥哥！醒醒！看着我！”

他试图靠近，但白千羽的攻击毫无章法却密集如雨，完全沉浸在自身的痛苦世界里。

几次险些被掌风扫中，白言冰瞅准一个空隙，不再犹豫，猛地欺近，不顾可能被击伤的风险，双臂张开，以一种坚定而不失温柔的力道，强行将狂乱的白千羽紧紧拥入怀中！

“千羽，别怕……哥哥在这里。”

他将下颌轻轻抵在她发顶，声音压得极低，却带着一种奇异的、直透人心的安稳力量，那是他魂骨境界自然散发出的沉凝气息，也是他此刻毫无保留的关切与守护之意。

“不管发生什么，哥哥都在。天塌下来，哥哥给你顶着。”

狂暴挣扎的白千羽，身体猛地一僵。

那熟悉的、令人安心的怀抱，那低沉而坚定的嗓音，如同穿透重重迷雾的阳光，一点点照进她混乱黑暗的心神。外界的流言、身世的颠覆、内心的恐惧……在这坚实温暖的怀抱面前，似乎暂时被隔开了。

“哥……哥哥？”她喃喃地、不确定地吐出两个字，赤红的眼眸中，疯狂之色渐渐褪去，取而代之的是巨大的迷茫、脆弱，和泫然欲泣的委屈。

“是我。”白言冰轻轻拍着她的背，如同儿时她做噩梦时那样，“没事了，没事了……”

白千羽紧绷的身体慢慢软了下来，将脸深深埋进兄长怀中，肩膀开始轻微地颤抖，无声的泪水迅速浸湿了白言冰胸前的衣襟。那是一种劫后余生般的依赖，也是一种情感堤坝崩溃后的宣泄。

白言冰就这样抱着她，任由她哭泣，心中充满了怜惜与愤怒。怜惜妹妹的无妄之灾，愤怒那泄露秘密、搅乱人心之徒。

他并不知道具体流言内容，但能让千羽如此失控，必定是触及了她最根本的恐惧。

不知过了多久，白千羽的抽泣渐渐停止，竟在白言冰怀中沉沉睡去，只是眉头依旧紧蹙，仿佛睡梦中也不得安宁。白言冰小心地将她抱起，放到一旁尚未损坏的软榻上，盖好锦被。他守在榻边，寸步不离。

又过了约莫一个时辰，白千羽睫毛颤动，缓缓醒来。眸中已恢复了清明，只是还残留着红肿与深深的疲惫。

她看着守在榻边的兄长，记忆如潮水般涌回，那些流言，那些痛苦……但奇异的是，此刻心中虽仍刺痛，却不再有那种足以让她崩溃的狂乱。

仿佛兄长那坚定的拥抱，为她筑起了一道临时的堤坝。

更让她惊讶的是，体内原本因走火入魔而紊乱不堪的气息，此刻竟平息下来，并且……似乎因祸得福？

在那种极端情绪冲击与后续兄长以沉凝气息安抚的过程中，她无意间冲破了某个关隘，此刻内视之下，骨骼隐现温润微光，虽然远不如白言冰的魂骨凝实璀璨，却更生机勃发，无疑已迈入了魂骨的门槛。

然而，修为的突破并未带来多少喜悦。一个更清晰、更尖锐的认知，随着神智的清醒，浮现在她心头：自己与眼前这个给予她无尽安全感的人并无血缘关系。

二十年来视作至亲、依赖敬爱的兄长，忽然变成了没有血缘关系的男子。

更进一步地，她甚至不属于这个家了。

一种陌生的、对于自己前途慌乱、恐惧的情绪，悄悄滋生。

她慌忙移开视线，不敢再与白言冰清澈关切的眼眸对视。

白言冰并未察觉妹妹心中这翻天覆地的变化，见她醒来且气息平稳，甚至略有进益，大大松了口气。“感觉如何？还有哪里不适？”

“……好多了，谢谢……。”白千羽低声道，声音有些沙哑。

她不知道现在该对白言冰采取什么称呼，于是生生地把刚到嘴边的“哥哥”两字咽了回去。

“到底发生了什么？为何会突然……”白言冰问道。

白千羽沉默片刻，苦涩地笑了笑，将听到的流言低声说了一遍。每说一句，白言冰的脸色就沉下一分，到最后，已是面罩寒霜，眼中厉芒闪动。

“荒谬！无稽之谈！千羽，你永远是我的妹妹，是父皇的女儿，是大周的公主！不要听信那些流言！”

白言冰斩钉截铁地说道，他本能地拒绝相信，或者说，拒绝接受这种会伤害到千羽的可能性。

白千羽看着他急于保护自己的样子，心中那丝异样的情愫又涌动了一下，但更多的是酸楚。她知道，哥哥这么说，是因为他爱护她。

可流言恐怕并非空穴来风。

安抚好白千羽，叮嘱宫人小心照料后，白言冰面色阴沉地走出暖阁。他需要立刻去见父皇，查清流言来源，严惩散布者！

然而，当他踏出宫门，沐浴在午后略显清冷的阳光下时，脚步却不由自主地慢了下来。方才暖阁中发生的一切——千羽的狂乱、脆弱，自己不顾一切的拥抱与守护，以及听到流言后的愤怒——种种画面与情绪在心头交织。

他忽然怔住了。

长久以来，他痴迷修炼，追求力量，最初是为了虚无缥缈的长生，后来是为了验证师尊传授的“道途”，在黑风岭是为了除暴安良，在皇城平叛是为了守护父皇和江山……

但直到此刻，当他看到千羽因身世流言而几乎自我毁灭时，一种前所未有的明悟，如同闪电般划过脑海。

他所追求的、所修炼的、所战斗的，本质上，不就是为了守护吗？守护这份让他感到温暖与责任的亲情，守护他在意的人不受伤害，守护这片给予他们安宁的土地！

长生也好，力量也罢，若失去了守护的对象，又有何意义？

“守护……我所珍爱的一切……”他低声重复，眼神从凌厉逐渐变得深邃而坚定。

心有所感，道有所应。就在这顿悟的刹那，他体内那早已魂骨圆满、内劲充盈的状态，忽然产生了奇异的共鸣。

骨骼深处那温润的金色光泽微微荡漾，骨髓中滋生的内劲不再仅仅满足于在经络中流转，而是自发地向骨骼更深处渗透、沉淀，仿佛要与之彻底融合，铸就更为不朽的根基。

他福至心灵，当即就在宫道旁寻了一处僻静角落，盘膝坐下，闭目调息。心神沉入体内，引导着那蜕变的内劲，按照冥冥中的感应，开始冲击那曾经以为虚无缥缈的“天骨”之境。

一切水到渠成。

夕阳西下，余晖洒落。静坐中的白言冰周身，渐渐泛起一层极其淡薄、却仿佛蕴含天地初开般古朴浩渺的微光。

他体内的骨骼，隐约传来大道纶音般的轻鸣，质地似乎发生了某种本质的升华，多了一份难以言喻的“天”之厚重与高远。

当最后一缕余晖没入宫墙，白言冰缓缓睁眼。眸中神光湛然，深邃如星空，却又平和如古井。

周身气息圆融内敛，再无半分锋芒毕露，却给人一种更加深不可测的感觉。

天骨，成。

然而，白言冰还未来得及体会这传说之境的神妙，更严峻的危机已如乌云般压城而来。

\vspace{2em}

关于白千羽身世的流言，已不再是宫闱秘闻，而是被摆上了朝堂，成了某些人手中攻击的利器。

关于安平公主白千羽身世的流言，如同附骨之疽，不仅在市井巷陌悄然滋生，更是在某些有心之人的推波助澜下，被堂而皇之地摆上了朝堂。这一日的早朝，气氛凝重得几乎能滴出水来。

以吏部尚书王焕、礼部侍郎郑儒林为首的一批“清流”大臣，手持祖宗法典，引经据典，言辞激烈。他们不再遮掩，直接将矛头指向了身处深宫的白千羽。

“陛下！”王焕须发皆张，声音洪亮，回荡在金殿之上，“臣闻，安平公主，实非陛下血脉，乃逆贼赵文渊与罪妃陈氏私通所出！此乃混淆天家血脉、玷污皇室清誉之滔天大罪！

“更有前日宫中异状，公主竟于深宫之内行止癫狂，毁物伤人，分明是修炼邪术、走火入魔之兆！此等妖女，岂可再容于皇室，玷污我大周国祚？按祖制，当以白绫赐死，以正视听，以儆效尤！”

“臣附议！”郑儒林紧接着出列，语气看似恳切，实则刀刀见血，“陛下仁厚，抚养其二十载，已是天恩浩荡。然国法祖制重于山，血脉正统不可乱。此女身负逆臣血脉，又心性不稳，易入魔道，留之恐成祸患。为江山社稷计，为皇室清誉计，请陛下忍痛割爱，明正典刑！”

一句句“妖女”、“私通”、“赐死”，如同淬毒的利箭，射向龙床之上那似乎苍老十岁的身形。

“放肆！”一声怒喝响起，白言冰终于无法忍耐，自武官队列踏出。他面色铁青，眼中燃烧着怒火，周身那天骨初成的沉凝气息不受控制地泄露出一丝，竟让离得近的几位文官感到一阵莫名的窒息心悸。

“千羽是我的亲妹，是父皇亲封的安平公主！什么流言蜚语，也敢拿到朝堂之上污蔑天家？王尚书、郑侍郎，尔等听信谗言，构陷公主，该当何罪！”

他的气势惊人，若在战场或江湖，确实足以震慑群雄；但这里是朝堂，是唇枪舌剑、规则博弈之地。白言冰空有撼山之力，却无辩倒众口之才。

他只会反复强调“兄妹”“公主”“构陷”，言辞直接甚至有些笨拙，在那些浸淫官场数十年、深谙言语机锋的文官面前，显得苍白无力。

“二殿下此言差矣！”王焕毫不退缩，反而挺直腰板，“臣等非是听信谗言，乃是据实直谏！公主身世，已有赵逆亲口供述为证，宫中走火入魔之事，更是多人目睹！此乃事实，岂是构陷？

“殿下武力超群，臣等佩服，然祖宗法度，岂能以武力凌驾？莫非殿下要效仿赵逆，以力压人，阻塞言路不成？”

这一顶大帽子扣下来，险恶至极。

“你……！”白言冰气结，拳头握得咯咯作响，却不知如何反驳。他感觉仿佛陷入了一张无形的大网，空有拔山之力，却无处施展。

他看向文官行列中的白承烨，后者此时正紧皱眉头，亦觉得如此紧逼实在不妥；但他很精明，只是装傻充愣、一言不发。

龙椅之上，白隐只觉得脑壳一阵阵抽痛。看着下方黑压压跪倒一片、口口声声“祖制”、“国法”的大臣，再看看急怒攻心却笨嘴拙舌的儿子，一股深深的无力感攫住了他。

养了二十年的女儿，乖巧懂事，医术高超，更救过自己的命，怎么可能没有感情？

可这满朝文武，铁板一块，以祖宗礼法相逼，他若强行维护，便是昏聩护短，势必动摇朝纲，寒了忠臣之心。治理国家，终究离不开这些人。

难道真要为了一个女儿（即便非亲生），与整个文官集团撕破脸，甚至大开杀戒？

那绝非明君所为。

争论持续了将近一个时辰，白言冰独木难支，文官们则步步紧逼，引经据典，慷慨激昂，仿佛今日不处死白千羽，明日就要亡国一般。

白隐被吵得头晕目眩，心烦意乱，最终忍无可忍，猛地一拍御案！

“够了！”他声音沙哑，带着疲惫与怒意，“此事……容朕三思！退朝！明日再议！”

说罢，也不顾群臣反应，直接拂袖而起，在内侍的高声“退朝”中，快步转入后殿，并严令侍卫将还想追进来劝谏的大臣全部“请”了出去。

回到养心殿，白隐颓然坐在椅上，揉着刺痛的太阳穴。

他感到自己仿佛被架在火上烤。

一边是难以割舍的亲情与对璟妃的愧疚，一边是稳固朝局的需要和祖制礼法的压力。

妥协？他做不到。强硬？似乎又力有未逮。

另一侧，白言冰正在和白承烨吵闹：

“皇兄为何不帮我说两句话？”

“那群老油条，父皇都说不动，我能说得过他们？难不成还能把他们全都砍了？

“那你倒是想点办法啊！”

“你都没办法，我能有什么办法？”

白隐的脑壳更疼了。

就在这时，他脑海中忽然闪过一个人影——那个总是神色平静、带着些许疏离与神秘的年轻人，叶牧谦。他能起死回生，能教出言冰这般鬼神莫测的弟子，或许……他真的有办法破此僵局？

“你们两个，不要吵了。承烨，你回东宫，今天不得出门。言冰，你去把千羽叫来。”

两个年轻人面面相觑，不知道皇上葫芦里卖的是什么药，但还是领旨去了。

很快，白隐、白言冰、白千羽三人未带任何仪仗，只做寻常打扮，悄然出宫，直奔叶牧谦在宫外的馆舍。

其时，叶牧谦正在院中槐树下，对着一盘残局皱眉苦思。他最近一直在推演养气境的细节，试图将“内劲”与更虚无的“元气”勾连起来，但这纯粹是理论空想，毫无实践基础。听到通报说皇帝私访，他连忙收敛心神，起身相迎。

馆舍简陋，白隐也顾不上讲究，屏退左右，仅留白言冰、白千羽兄妹。随即，便将朝堂上那令人窒息的压力与自己的两难处境，向叶牧谦和盘托出。

“叶先生，朕……朕实在是没有办法了。”白隐脸上满是疲惫与无奈，“千羽是朕看着长大的孩子，虽然近日才知道不是朕亲生的，但依然视如己出。可那群大臣……唉。朕既不能寒了忠臣之心，乱了朝纲，又绝不能伤害千羽。

“先生学究天人，神通广大，可有良策教我？”

叶牧谦静静听完，手指无意识地在石桌上轻轻敲击。虽然他不是权谋之士，但至少来自信息爆炸的时代，看过的宫斗剧、权谋小说也不知凡几了。结合当前情况，一个大胆的、带有表演性质的计划，逐渐在脑中成型。

他抬起眼，目光扫过满脸期待的皇帝，神情紧绷的白言冰，以及眼中带着最后一丝希冀的白千羽，缓缓开口：“陛下，此事的关键，在于‘势’。文官集团挟祖制礼法以造势，逼陛下就范。若要破局，须以更强的势，压过他们。”

“更强的势？”白隐疑惑。

“不错。”叶牧谦点头，“明日朝会，陛下可假意妥协，召见千羽，当庭宣布其‘罪状’，并做出赐死之态。”

“什么？！这绝对不行！”白言冰失声，白千羽也脸色一白。

“稍安勿躁。”叶牧谦摆手，继续道，“届时，我会适时出现。”他顿了顿，眼中闪过一丝难以察觉的微光。

这计划固然是为了解围，但何尝不是一次绝佳的展示机会？

“我会展示一些超乎他们理解的力量。”叶牧谦语气平静，却带着一种莫名的分量，“以此震慑群臣，让他们明白，有些存在，非世俗礼法所能约束。”

他接着补充：“但是，如果我只是这样救下千羽，那么皇室必然威信扫地，这不利于朝堂的稳定。故而，象征性的惩罚，自然要有；我对贵国礼法祖制一窍不通，这还需陛下与言冰和千羽商议决定。

“如此，既保全了千羽性命，也全了陛下慈父之心，更给了文官集团一个台阶下，同时还维护了皇室的威严与礼法。”

白隐听罢，眼中光芒闪烁，仔细思量。此计虽险，却似乎是在绝境中唯一能兼顾各方的办法。

他看向自己的一双儿女，似在征求他们的意见。

白千羽适时的盈盈下拜，声音哽咽却坚定：“千羽……愿听从安排。只要能不再让父皇和兄长为难。只是……我希望自此搬离皇宫，随师尊清修。千羽不愿再接近那危险的朝廷。”

“不行，你……”白言冰眉头紧锁，他厌恶这种算计。

“若千羽欲随我清修，也可。”叶牧谦回答道。

在白隐和叶牧谦看来，白千羽果然冰雪聪明：离开皇宫，确实能让她远离风暴中心。

但此时，她内心的实际想法，却是觉得无法再面对那对自己倾注了一切守护之意的白言冰——只是在场无人能理解她的这份心罢了。

为了妹妹的安危，白言冰咬牙点头：“但凭父皇和师尊做主。”

叶牧谦又补充道：“既然明日有我能压过群臣，那陛下自然可以做些激进的变革。我心知陛下见文官集团铁板一块，内心不爽；何不借此机会，暗中分化他们？”

白隐的眼睛猛地一亮。对啊，他怎么没想到？这位先生真是老成谋国之辈。

“朕知道了。”白隐点了点头。“明日，朕自然会配合先生演出这场大戏。”

计划就此敲定。

又略略商议了一会，白氏三人怀揣着复杂的心情离去。只是不知道，明日朝堂之上的这场表演，究竟会如何。

\vspace{2em}

次日，大朝会。气氛比昨日更加肃杀。

白千羽一身素衣，跪在殿中；文武百官分列两旁，许多人都目光闪烁，等待着皇帝最终的裁决。

白隐高坐龙椅，面色沉痛，缓缓开口，陈述了白千羽身世存疑、宫中行止失常等“罪状”，语气沉重，仿佛经过了极大的挣扎。最后，他闭了闭眼，挥手下令：“安平公主白千羽，虽朕抚养多年，然祖制不可违，国法不可废……赐白绫。”

“父皇！”白言冰目眦欲裂，上前一步，却被两名殿前侍卫微微拦阻。换在往常，这两名侍卫怎么可能会拦白言冰的行动？不过是预先得了点拨，跟着逢场作戏罢了。

白千羽闻言身体微颤，却依计深深叩首，声音清晰：“罪女……领旨。”

一名手托银盘、上覆白绫的内侍，低着头，迈着细碎的步子，从侧殿走出，走向白千羽。那刺目的白色，在庄严的金殿上，显得格外残忍。

殿中鸦雀无声，许多大臣低下头，不忍再看；王焕、郑儒林等人，虽然以笏板掩住眼睛，嘴角则难以抑制地露出一丝得意。内侍走到白千羽身前，即将展开白绫，就在那千钧一发之际——

“且慢！”

一个清朗平静，却仿佛蕴含着奇异穿透力的声音，自殿外响起。这声音不高，却清晰地传入每个人耳中，甚至盖过了殿内所有的细微声响。

众人愕然望去。

只见殿门处的阳光被一道身影遮挡。叶牧谦穿着一身简单的青色布袍，负手而立，缓缓步入金殿。他步伐从容，目光平静地扫过殿中众人，最后落在白千羽身上。

守卫殿门的御林军精锐，竟无一人上前阻拦，甚至在他目光扫过时，下意识地后退了半步，仿佛被某种无形的气势所慑。文武百官中，更无一人敢在此刻出声呵斥这个擅闯金殿的布衣青年。整个朝堂，竟因他一人的到来，陷入了一种诡异的静默。

叶牧谦大剌剌的走到大殿中央，看向龙椅上的白隐，又环视群臣，目光终于落到那银盘之上。随即，眉头一皱，声音也冷了下来：“我的宝贝徒儿，谁敢动她？”

这怒意，三分是装出来的，七分是真的。

话音未落，他抬起右手，食指与中指并拢，看似随意地朝着身旁一根需两人合抱的蟠龙金柱，隔空一划！

没有惊天动地的声响，没有光华四射的异象。

但所有人都清晰地看到，随着他手指划过的轨迹，那坚硬无比、象征着皇室威严的蟠龙金柱表面，仿佛被一柄无形却无比锋利的巨刃划过，发出“嗤”一声轻响。

一道深达三寸、光滑如镜的笔直划痕，凭空出现在金柱之上！切痕边缘，石粉簌簌落下。

隔空数丈，裂金石！

整个金殿，死一般的寂静。所有声音都被抽空了，只剩下众人粗重或屏住的呼吸声。王焕手中的笏板“啪嗒”一声掉在地上，郑儒林脸上的得色瞬间冻结，化为难以置信的惊恐。就连知晓计划的白隐和白言冰，此刻亲眼见到这超出理解的一幕，心脏也猛地一缩。

这绝非武功！这是……仙术？妖法？还是别的什么？

叶牧谦收回手指，仿佛只是做了一件微不足道的小事。

他再次看向白隐，语气平淡，却带着不容置疑的意味：“陛下，千羽是我弟子。她若有错，自有我这师尊管教。只是……世俗律法，怕是管不到我门中之事吧？”

白隐恰到好处地露出了惊惧之色，仿佛也被叶牧谦这一手吓住了，连忙顺势道：“仙……仙师息怒！是朕……是朕考虑不周！”他擦了擦额头上并不存在的冷汗，快速道：“安平公主白千羽，即日起，削去公主封号，贬为安平郡主！废国姓白，改随其母姓陈！责令其即刻搬离皇宫，随叶仙师清修，未经传召，不得擅入宫闱！”

一连串的处罚决定快速宣布，从赐死到贬斥流放，转折之快，令人目不暇接。

跪在地上的白千羽——现在应该叫陈千羽了——再次叩首，声音平静无波：“罪女陈千羽，领旨谢恩。”

叶牧谦不再多言，对白隐微微颔首，便转身，带着陈千羽，在满朝文武呆若木鸡的注视下，从容步出金殿。殿外，白言冰早已安排好的简朴马车等候着。两人登上马车，车夫一挥鞭，马车便辘辘驶离了皇宫，驶离了这个刚刚经历了一场无声惊雷的是非之地。

直到马车消失在宫门之外，金殿内的寂静才被打破，随即响起一片压抑的、难以置信的嗡嗡议论声。但再也没有人敢高声提出异议。那道深达三寸的柱上刻痕，如同烙铁，深深印在了每个人的心头。那是一种超越他们认知范畴的力量，足以粉碎一切所谓的礼法、祖制和阴谋。

白隐强压住心中的波澜，看着下方神色各异的臣子，趁势宣布了另一项决定：“另，宰相之位空悬，国事不可久旷。即日起，由太子白承烨暂领相职，参知政事，辅佐朕处理朝务。”

白承烨正在下方装傻充愣。忽然听到父皇宣布如此决定，内心狂喜，却强行压住，只是迈出一步，恭敬道：“儿臣遵旨。”

王焕等人张了张嘴，看着皇帝那不容置疑的脸色，又瞥了一眼那根柱子上的刻痕，终究将反对的话咽了回去。不少臣子内心已经开始盘算，要怎么开始巴结白承烨了。

经此一事，白隐对叶牧谦的敬畏，达到了前所未有的高度。那不仅仅是救命之恩、教导之德，更是一种对未知而强大力量的深深忌惮与倚重。

马车驶入叶牧谦所居馆舍的后院。车门打开，陈千羽在叶牧谦的示意下先下车回房休息。当车帘落下，车厢内只剩下叶牧谦一人时，他挺直的腰背瞬间佝偻下去，脸色以肉眼可见的速度变得苍白如纸。

“噗——”

一口暗红色的鲜血，毫无预兆地喷溅在车厢地板上。紧接着，四肢百骸仿佛被无数细针攒刺，酸软麻木；尤其是并指划出的右臂，此刻又像是被重锤反复敲打般的剧痛，几乎抬不起来。一股深入骨髓的疲惫和空虚感，如同潮水般将他淹没。

他强撑着，用袖子擦去嘴角血迹，踉跄着下车，回到自己房中，立刻盘膝坐下，试图调息。然而，体内那股平日因“反哺”而存在、今日被他强行调动起来模仿“隔空剑气”的力量，此刻狂暴紊乱，反噬自身，根本不受控制。

“原来……强行动用，代价如此之大……”叶牧谦额头冷汗涔涔，心中恍然。

他之前只是理论推演，从未真正尝试过将这种力量外放用于攻击或展示。今日为了震慑朝堂，强行催动，果然遭到了剧烈的反噬。这让他彻底明白，自己这身莫名得来的力量，绝非可以随意挥霍的玩具。每一次动用，尤其是超出当前身体承受范围的动用，都可能带来严重的后果。

隔壁房间，陈千羽换上素衣，默默收拾出一间临街的静室，挂上一块简单的木牌，上书“安平医舍”。她决定在此行医，仅收取药材成本及十分之一的诊费，不为牟利，只为在治病救人的实践中，继续深化对医道、对生机、对“元气”的理解。这是她选择的，远离风暴后平静而充实的生活方式。

“安平医舍”的木牌刚刚挂起，白言冰也很快赶来。

陈千羽见到他，眼睛微微睁大了一瞬。

“殿下怎么来了？”

她刻意地调整了她的措辞。

白言冰笑道：“朝廷可太虚伪了，于我本性不合。我已经跟父皇说了，此后除非国家有倾覆之危或父皇命令，否则我可不再参与日常朝政。”

随即又说：“咱俩私下里没有遵守皇室礼节的必要，还是用你我相称吧。这个‘殿下’，叫得太生分了，真难听。”

听闻此语，陈千羽嘴角弯了弯，心中泛起一股暖流。她偏过头去，调整了一下表情，又问：“那，师兄，东宫怎么说？”

白言冰忽然觉得这句“师兄”叫得太调皮了，似乎小猫在他心上抓挠。

“承烨肯定不可能强留我啊。私下归私下，在朝廷上他只会嫌我帮倒忙。”

随即又面向刚刚在车上调息完成、走下车的叶牧谦。

“师尊，以后我也就住这了。吃喝方面，宫里会定时派人送来，师尊就不必担心了。”

叶牧谦心底稍微惊奇了一瞬，不过很快便为更大的了然所取代。

“嗯。”

自此，白言冰与陈千羽日常便泡在叶牧谦这方小小的天地中。叶牧谦在休养恢复的同时，也开始正式引导两人进行养气境的初步修炼，教导他们如何感知、温养那虚无缥缈的元气，为真正踏入下一境界打下基础。

当然，在所有人意料之中的是，自那日金殿一指之后，无论是原本对陈千羽身份耿耿于怀的官员，还是可能存在的赵党余孽，甚至是京城其他势力，竟无一人敢来这间看似普通的馆舍闹事。

很简单，打不过。

\vspace{2em}

皇城的风波似乎随着陈千羽的离开而暂时平息，叶牧谦那间位于京城僻静处的小院，成了师徒三人难得的安宁港湾。

有不少百姓因陈千羽治好了积年的疾病而感恩戴德，甚至有人送来一袋玉米……白言冰、陈千羽虽不能说是娇生惯养，但也是皇家血脉；宫内每个月也会送来粮食，他们也不缺吃的。如果是玉米面倒也无所谓了，要玉米有何用，难不成要自己磨成碴子？但那老头子强塞，陈千羽也不好拒绝，于是这玉米只能一直在墙角堆着。

日子在规律的修炼与学习中悄然流逝。院角的药圃被陈千羽打理得生机勃勃，白言冰则在树下开辟了一小块空地，每日勤练不辍，打磨着天骨境界那愈发圆融沉凝的力量。不忙时，陈千羽的眼睛则常常盯着白言冰练剑的姿态微微出神。

叶牧谦本人则完全一副游戏人间的高人姿态，对陈千羽这微妙的变化看破不点破。

虽然他强行动用养气境造成的反噬消退，但那种深入骨髓的疲惫和动用超越界限力量后的空虚感，让他心有余悸。他更加确信了自己之前的猜想：这套“问道途”的力量，对他而言并不是一些可以随意取用的宝藏。

这一日，阳光透过槐树叶的缝隙，洒下斑驳的光影。陈千羽在晾晒完一批药材后，看着手中一本不知从哪个书摊淘来的、描绘着仙人炼丹的粗劣话本，忽然抬起头，眼中闪着好奇与探究的光芒，望向正在石桌旁对着自己那堆鬼画符一般的推演图发呆的叶牧谦。

“师尊，”她轻声开口，指了指话本上那被描绘得金光闪闪、云雾缭绕的丹炉，“这书中常说的‘丹药’，究竟是何物？与我们平日煎煮的汤药，有何不同？”

叶牧谦闻言，心中顿时一咯噔。

丹药？这触及他的知识盲区了！

他前世就是个写网文的，对于炼丹的认知仅限于小说里“天材地宝、文武火候、丹成几转”之类的套路词汇，具体怎么操作，那是一窍不通。但看着徒弟那求知若渴的眼神，他只能硬着头皮，努力让自己的语气听起来高深莫测。

“丹药么……”他沉吟着，搜刮着肚里那点可怜的存货，“其本质，亦是萃取药材精华，调和阴阳，只不过形式不同。汤药乃是以水为媒，化药力于汤液，便于吸收，但不易保存携带。而丹药，则是以特殊法门，将药力高度浓缩凝聚成固态丸剂，胜在便携、效力持久，且对于修士而言更易于控制一些猛烈药性的释放节奏。但是对于凡人，这种丹药反而容易因为药力过猛、无法承受药力对四肢百骸的冲击，从而伤了身体。”

他尽量说得玄乎，把浓缩和便携作为核心区别，这至少符合现代人对药丸和汤药的朴素认知。

陈千羽若有所思地点点头，眼中却燃起了更旺的探究之火。“浓缩……便携……师尊，那这‘特殊法门’，究竟是如何将药力浓缩成丹的呢？弟子……可以试试吗？”

叶牧谦本想直接拒绝，但转念一想，让她去碰碰壁也好，或许能让她知难而退，便模棱两可地道：“炼丹一道，博大精深，涉及火候、药材君臣佐使、甚至天地时令，非一日之功。你若感兴趣，可先从最基础的尝试起，但切记，莫要强求，更不可胡乱服用不明之物。”

更何况，他在这个世界，可不敢确定陈千羽到底会不会碰壁了：万一真做出来了些能用的丹药呢？

他这敷衍的鼓励，在陈千羽听来却成了许可。这位在医道上向来执着的姑娘，立刻将研究炼丹列为了新的课题。她就地取材，用院里煎药的陶罐和小炭炉开始了最初的尝试。

她的思路直接而朴素：既然丹药是浓缩的汤药，那我把汤药里的水熬干不就行了？

于是，她精心调配了一剂最熟悉的补气养血汤药，照常煎煮，待药汁浓缩后，并不倒出，而是继续用文火小心熬煮，试图将最后的水分蒸发，得到“药膏”乃至“药丹”。

结果可想而知。水分即将熬干之际，罐底只剩下一层焦黑粘稠、散发着糊味的药渣，不仅原本清亮的药汁变得漆黑，那股焦糊气更是刺鼻。陈千羽用银勺小心挑起一点，放在鼻尖轻嗅，又舔了舔，眉头紧锁——药性全无，只剩下焦苦之味，甚至可能产生了有害物质。

第一次尝试，以失败告终。

叶牧谦踱步过来，看着罐底那团焦黑，心中暗松一口气，表面却不动声色，借坡下驴道：“凡品药材，大多草木之属，其内蕴药性精华，往往不耐持久高温炙烤。你看，水分蒸发殆尽之时，亦是药性挥发或变质之际。修士所用炼丹之材，多为天地灵物，本身便能耐受烈火焚烧，甚至需以真火、地火长时间淬炼，方能逼出精华。以凡火、凡药炼丹，火候极难掌控，稍有不慎，便是这般结果。”

他这番话半是真知灼见，半是胡诌。

然而，陈千羽的倔强此刻被彻底激发了出来。她不肯轻易认输。“火候难控……那就控制到极致！”她抿着唇，眼中闪烁着不服输的光芒。

接下来的日子，小院的角落成了陈千羽的“丹房”。她弄来更精细的炭火，甚至尝试用不同木材烧制的木炭，以控制火力的温和与持久。她找来薄铁片置于炭火上，再放陶罐，试图让受热更均匀。她将药材研磨成极细的粉末，或先以酒浸，或先以蜜炼，尝试各种预处理方法。她屏息凝神，全神贯注地盯着陶罐内的变化，用银针、用观色、用闻味，试图捕捉那稍纵即逝的“成丹”时机。

但无论她如何精微地操控火候，变换方法，总是在药汁即将彻底干涸、似乎要凝结成型的最后关头，功亏一篑。不是突然焦糊，就是药力散逸，得到的总是一团性状不定的焦糊物或松散的药粉，根本无法凝聚成丸。

屡次失败，并未让她气馁，反而促使她更深入地思考。

一日，她正对着又一次失败的药渣发呆，脑海中忽然灵光一闪，如同黑暗中被闪电照亮！

“水……汤药中的水，其作用究竟是什么？”

水，不仅仅是将药材中的药力逼出，更是作为药力精华的载体！药材中的有效成分溶于水，随汤液进入人体，方能发挥作用。若直接将水熬干，作为载体的水消失了，那些被提取出的药力精华，便失去了依附之物，要么随水汽挥发，要么在高温下分解变质，自然无法成丹！

想通了这一点，她豁然开朗！问题的关键，不在于去掉水，而在于找到一个能替代水、且能在高温下保持稳定、并能最终承载药力成型的新“载体”！

她的目光扫过院中，最终落在了墙角那袋晒干的玉米上，忽然想到：煮过玉米的水，会变得粘稠，那东西厨子叫“淀粉”，倒是可以用来勾芡……一个大胆的创新想法在她脑中成形。

她立刻行动起来。取来玉米，加水熬煮，得到一锅略显浑浊粘稠的淀粉水。然后，她不再直接用清水煎药，而是用这过滤后的淀粉水作为媒介，投入按比例配好的药材，开始小心熬煮。

过程依旧需要极其精细的火候控制。淀粉水比清水更容易糊底，温度高了、时间久了，就会焦糊。陈千羽几乎将全部心神都投入其中，小心翼翼地调节着炭火，观察着药液的变化。

这一次，情况有所不同。随着水分慢慢蒸发，罐中的药液并未直接变焦，而是逐渐变得极其粘稠，颜色也转为深褐。药力似乎被淀粉形成的胶体网络牢牢吸附、包裹。在水分将尽未尽的某个微妙时刻，陈千羽果断撤去火源，用浸过凉水的布包裹住滚烫的罐子，让其快速降温。

待罐子冷却，她用竹刀小心刮取罐壁和内底。不再是焦黑的药渣，而是数颗大小不均、表面略显粗糙、呈深褐色、触手微硬却带着弹性的……丹丸！虽然卖相不佳，甚至能闻到淡淡的玉米味混杂药香，但它们确确实实是成型的丹丸！

“成……成功了？”陈千羽捏起一颗尚带余温的丹丸，难以置信。她小心刮下一点粉末尝了尝，眼中顿时爆发出惊喜的光芒——药味纯正，虽然比汤药似乎淡了一些，但确确实实是那剂补气养血方的味道！而且，似乎更温和，更易于保存了！

她捧着这几颗堪称简陋的“丹药”，如同捧着珍宝，快步跑到叶牧谦面前。

叶牧谦正对着自己画的元气运行图打瞌睡，被陈千羽唤醒，看到那几颗其貌不扬的丹丸，又听了她的叙述，内心顿时掀起惊涛骇浪！

成了？居然真的用凡火、凡药、甚至用玉米淀粉水这种匪夷所思的“载体”，炼出了成型的丹药？

虽然这丹药品质一看就极其低劣，远不能和传说中金光闪闪的灵丹相比，但这毕竟是“丹”！这意味着，陈千羽不仅理解了他随口胡诌的浓缩概念，更凭借自身的医学知识和不懈实验，硬生生在一条看似不通的路上，开辟出了一个可行的方向！

他强压下心中的惊诧，面上露出赞许之色，接过丹丸仔细看了看，又嗅了嗅，点头道：“不错。能以凡品药物、凡火成功凝丹，虽借了淀粉作为形骸依附药力，但这份对火候精微的掌控，已堪称炉火纯青。只是……”他话锋一转，将丹丸放下，“受制于药材本身品质低劣，此丹效力平平，且杂质颇多，丹毒未清，算不上真正的‘灵丹’。真正的上品丹药，当以精纯药力自行凝聚，浑然天成，无需这等外物辅料参与调和。”

当然这话大多还是胡诌。

陈千羽听得认真，眼中光芒更盛，似乎又有了新的目标。她忽然抬起头，带着些许俏皮和探究，问道：“师尊，那弟子这初学之作，比之当年初学炼丹的您……如何？”

叶牧谦被问得一愣，随即心中苦笑。

比当年的我？我当年连个药罐子都没摸过！

但他脸上丝毫不敢露怯，只能干咳一声，故作高深地捋了捋并不存在的胡须，扯谎道：“嗯……单论这初次成丹的品相与控制力，比初学时的为师，却是强上不少！”他看到陈千羽眼中闪过一丝得意，连忙又补充道：“不过，‘闻道有先后，术业有专攻’。你本就精研医理药性，是为师之长。而在其他方面，你师兄或许又有其优势。切不可因此自满。”一番话，既抬高了徒弟，又模糊了焦点，还顺便夸了一下旁边的白言冰，端的是滴水不漏。

陈千羽抿嘴一笑，不再追问，珍而重之地收起那几颗初成的“淀粉丹”，心中充满了继续钻研的动力。

然而，就在陈千羽成功炼制出那几颗简陋丹药的当晚，叶牧谦在独自静坐、酝酿睡意时，忽然感到脑海中一阵清凉，许多关于火候控制、药性萃取、君臣佐使在凝丹中的变化、淀粉作为载体的优缺点等纷繁复杂、却又条理清晰的知识与感悟，凭空涌现！仿佛他亲身经历了无数次失败的尝试，最终总结出了陈千羽那套“淀粉水炼丹法”的所有关窍与细节！

他顿时睡意全无。猛地睁开眼，眼中满是不可思议。

“这……这是……千羽炼丹的经验和知识？难道说，不只是修为境界，连她自行研究领悟出的技艺，也会‘反哺’给我？”

这个发现让他心脏狂跳。如果真是这样，那这“问道途”的奥秘，远比他想象的更加诡异和强大。

而强行越阶造成反噬这一点，与如此巨大的发现相比，似乎也不是什么大不了的问题了。

嗯，这么说归这么说，对叶牧谦本人而言，他是不想再吐第二次血的。

不过很多时候由不得他。

\vspace{2em}

就在次日，一直在潜心温养元气、尝试迈入养气境的白言冰，终于取得了突破。

院中槐树下，他盘膝而坐，呼吸逐渐变得绵长深远，仿佛与周围环境融为一体。体内那天骨境界滋生的内劲，在某种玄妙的引导下，开始与呼吸、与周身毛孔感知到的某种极淡的、无所不在的微薄能量——也就是被叶牧谦称为“元气”的那东西——尝试勾连、融合。

不知过了多久，白言冰身体微微一震，周身气息陡然一变，一股温润而充满生机的气流自他丹田处升起，沿着某种特定的路径，缓缓在他体内运行了一个完整的循环——第一周天，完成！

——吾善养吾浩然正气（《孟子·公孙丑上》）！

养气境，初成！

就在白言冰体内元气成功转过第一周天的刹那，隔壁房中正在消化新得炼丹知识的叶牧谦，猛地感到自己小腹丹田之下，仿佛有什么东西轰然炸开！

是一种豁然开朗、空间拓展的奇异感觉！

他连忙闭目内视，震惊地“看”到，在自己丹田下方，原本混沌一片的区域，不知何时竟然自行开辟出了一片朦胧的、泛着微光的空间！那空间似虚似实，仿佛能容纳无尽的气息——正是他胡编乱造的养气境核心标志之一，气海！

而且，与之前强行催动养气能力后遭受反噬不同，此刻他心念微动，尝试调动一丝力量，感觉那新开辟的气海中，一股温顺而充盈的暖流随之响应，运转自如，再无半点滞涩，使用后也无半分亏虚和痛苦！

叶牧谦呆坐了半晌，一个清晰的认知逐渐在脑海中成型：

他胡编的这套理论框架，其中未被徒弟走通的部分，若强行使用，便会遭到世界或体系本身的反噬，代价惨重。但是，一旦徒弟凭借自身努力，真正踏足了某个境界，领悟了某种技艺，那么这条路就被真正走通了，变成了现实；而叶牧谦作为这个体系的源头，便能自然而然地同步获得相应的能力，且再无使用障碍！

想通了这一点，他心中既感震撼，又生出一种难以言喻的明悟。叶牧谦这套道途，与其说是他传授给徒弟的功法或说是修炼法门，不如说是一个由他提出设想、由徒弟去探索实现、最终双方共同受益的奇异共生体系！

白言冰的第一反应，自然是去找叶牧谦汇报。

他如今养气初成，感知更为敏锐，下意识地便想探查一下深不可测的师尊。当他小心翼翼地探出一丝元气感知叶牧谦时，却仿佛泥牛入海，只感到师尊丹田之下那片气海，浩瀚深邃、无边无涯，以他初入养气的微末修为，根本探不到底！

白言冰心中骇然，对叶牧谦的敬畏瞬间达到了新的高度。

师尊果然早已远远走在了前面，其境界如渊如海，深不可测！

敬畏之余，白言冰又想起自己如今虽力量增长，但剑法招式依旧主要靠自悟和实战总结，缺乏系统传承，便恭敬请教：“师尊，弟子如今已初步养气，不知师尊可否传授一些精妙的剑法招式？也好让弟子在运用元气时，能有更恰当的施展之法。”

叶牧谦闻言，心中又是一阵发虚。

剑法？他一个现代文科生，体育课都没怎么认真上过，哪里懂什么精妙剑法？

太极剑倒是知道个名字，但具体招式？早就还给公园老大爷了。

他只能再次端起师尊的架子，故作深沉道：“剑法之道，存乎一心。招式是死的，人是活的。你已初步养气，对自身力量、对周围环境的感知远超以往。真正的剑法，当从实战中体悟，从自然中借鉴，与你自身的‘道’相结合。为师若给你固定套路，反而可能限制你的发展。你不妨多回想战斗经历，观摩风雨雷电、山川走势，从中领悟属于自己的剑意。”

一番话说得云山雾罩，核心思想就一个：自己琢磨去！

白言冰听了，虽然觉得有道理，但眼神中还是不免流露出一丝失望和淡淡的不满。他总觉得师尊似乎在某些非常具体的“术”的层面，格外吝啬或含糊。但他不敢质疑，只能躬身应是，心中暗下决心，定要自己悟出一套配得上如今修为的剑法来。

\vspace{2em}

皇城金殿一指刻柱的余波，远比叶牧谦想象的更为绵长深远。

白言冰一人一剑独挡数千叛军的英姿，陈千羽猛药救父、起死回生的传奇，连同他们背后那位神秘莫测、被二皇子与公主共同尊为师尊的布衣青年——叶牧谦的故事，如同插上了翅膀，以惊人的速度在京城乃至更远的地方流传开来。

说书先生们将之编成跌宕起伏的话本，街头巷尾的百姓津津乐道，茶余饭后，“无上仙师”、“护国仙师”的名头不胫而走，叶牧谦虽深居简出，其声望却在无形中被推至了一个凡人难以企及的高度。

随之而来的，是无数炽热乃至狂热的目光。

京城中诸多勋贵世家、豪商巨贾的纨绔子弟，平日里鲜衣怒马、斗鸡走狗，此刻听闻世间竟有如此“真仙”在世，能授人以超凡脱俗、甚至可叩长生、登极巅的法门，哪里还坐得住？

求仙访道、渴望一步登天的欲念，瞬间压过了所有。

他们不再满足于听那些说书人对叶牧谦的吹捧，开始携带着成箱的金银珠宝、珍奇古玩、甚至家传的所谓“灵物”，络绎不绝地涌向叶牧谦那间位于僻静处的馆舍，意图叩开仙门，拜入这位“无上仙师”座下。

然而，叶牧谦对此的反应，却是一以贯之的冷淡与拒绝。院门终日紧闭，对于门外喧嚣的求见声、丰厚的献礼，他只让白言冰、陈千羽，或偶尔前来馆舍中帮忙的老仆传出一句话：“师尊清修，不见外客。”

这并非故作姿态：叶牧谦透过门缝或是在白言冰的转述中，早已将门外那些求师者的形貌心性看了个大概。多数人眼神浮躁，气息虚浮，言谈间不是炫耀家世便是许以重利，更有甚者，身后跟着的恶仆依旧带着欺压良善的跋扈之气。

这些人，不过是慕虚名、求捷径，欲借“仙缘”锦上添花，甚至满足更不堪的私欲，何曾有半分问道求真的诚心与坚韧？此其一。

其二，叶牧谦深知自己根底。他所凭依的道途诡异莫测，与徒弟命运深度绑定，且蕴藏风险。收徒绝非儿戏，现在肯定不是开山立派、广纳门徒的时候。这些纨绔背后往往牵扯着盘根错节的势力，一旦沾染，无穷无尽的麻烦便会接踵而至。他只想守着这方小院，安静地引导白言冰和陈千羽走下去，慢慢弄清自身与这个世界的奥秘，绝不愿卷入更多的世俗纷争与权力倾轧。

纵然叶牧谦有这样那样的考量，但他的闭门谢客，在那些心高气傲的纨绔子弟看来，却是莫大的羞辱与轻蔑。

“不过一个装神弄鬼的江湖术士，真当自己是仙人了？”

“定是嫌礼不够重！或是故作清高，待价而沽！”

“听说他来自十万群山之外？莫非仙缘就在那群山之中？”

“是了！定是如此！他能在那里得道，我们为何不能？与其在此受他冷眼，不如我们自己去寻那仙缘！”

\vspace{2em}

流言与愤懑交织，催生出了愚蠢的勇气。

一部分最为胆大、又对“仙道”执念最深的纨绔，竟然私下串联起来。他们不顾家人劝阻（有些甚至瞒着家里），纠集了一帮重金聘来的、武功还算不错的护卫，携带了大量他们认为“探险”所需的物资，在一个雾气朦胧的清晨，悄然离开京城，向着西方那连绵无尽、被当地人敬畏地称为“十万群山”的苍茫山脉进发。

他们天真地以为，只要翻越这片传说中的险地，就能抵达仙师来的地方，找到真正的长生秘境、无上仙法。

当然，等待他们的，自然不是仙境。

十万群山，因其山岭众多、层峦叠嶂而被冠以“十万”之名。最末一峰——云海腾涌、山势奇秀、路途危险的绝念峰，也就是白言冰、陈千羽遇到叶牧谦的地方——是整个十万群山环境最宜居或最“正常”的。

越是深入十万群山，环境越是凶险：终年不散的毒瘴之气弥漫在深谷幽壑，吸入少许便可能令人头晕目眩，严重者甚至脏腑溃烂；密林深处，潜藏着无数凶猛的异兽——被称为“妖兽”者，虽能看出与常见飞禽走兽有些相似之处，但力量、速度等都比常规的野兽要强得多：有些体型大如房屋，力能裂石，有些则小巧迅捷，齿爪带毒，防不胜防；更有诡异的地形，看似平坦的草地可能是吞噬一切的沼泽，宁静的水潭下或许藏着致命的水怪。

这群养尊处优的公子哥和他们的护卫，甫一进入群山外围，便吃足了苦头。昂贵的绸缎衣衫被荆棘刮得破烂，沉重的金银珠宝成了逃命的累赘。毒瘴让他们咳嗽不止，脸色发青；妖兽的袭击接二连三，护卫们拼死抵抗，伤亡惨重。他们被迫放弃大部分行李，像没头苍蝇一样在深山老林中东躲西藏，很快便迷失了方向。

疲惫、恐惧、疾病交加，最初的雄心壮志早已被残酷的现实碾得粉碎，只剩下求生的本能和对归途的渴望；但已经走到这一步了，沉没成本太多、回去还太丢脸了，于是只能在这里继续吊着一口气挣扎：万一真能遇到仙人呢……

\vspace{2em}

群山另一侧，那被某些存在称为“内域”的世界里，一场追杀刚刚告一段落。

邪修“血魔”，修炼的是以生灵鲜血祭炼邪气、滋养己身的歹毒功法，因其手段残忍、杀戮无度，在内域一片区域恶名昭彰，终于惹来了正道修士的刺杀。一场激战，血魔虽仗着功法诡谲重伤逃遁，潜入十万群山，但其本源已损，实力大跌。

追杀他的修士追至群山边缘，感应着山中驳杂混乱的气息和浓郁的瘴气，判断血魔重伤之躯难以在如此恶劣环境中长久支撑，更别提翻越这广阔险峻的群山了；而且十万群山对面，那可是荒域啊，灵气不显、更无修士；修士就算是到了那里，也会很快因为外界灵气浓度不足，灵枢中的天地灵气发生泄漏，从而跌落成凡人。她经过思考与判断，最终决定在内域这一侧守株待兔，只等血魔熬不住退回时，一举擒杀。

血魔的思维几乎一致。他也不是糊涂蛋，深知以现在这种情况，返回内域必死无疑。感应到身后隐约的封锁气息，他把心一横，竟不顾一切地向着群山更深处、那天地灵气明显变得稀薄贫瘠的方向——也就是“荒域”——亡命逃窜。他虽倚重灵气，但短时间内，利用邪法以大量生灵精血强行祭炼出的邪气，亦可替代灵气，支撑他行动甚至恢复部分实力。

而荒域虽传闻灵气稀薄，但……总有生灵吧？

\vspace{2em}

于是，一个重伤急需“补给”的邪魔，与一群迷失山野、筋疲力尽的猎物，在十万群山深处一片雾气弥漫的山谷中，“幸运”地相遇了。

\vspace{2em}

对于几乎弹尽粮绝、被妖兽和瘴气折磨得奄奄一息的纨绔与护卫们来说，突然出现的血魔——尽管其形貌因伤势和功法显得极为狰狞可怖——起初甚至被他们误认为是山中的猎户或同样迷路的旅人，几乎要喜极而泣。

而对于血魔而言，眼前这群气血虽然不算旺盛、但数量可观且毫无反抗之力的人类，简直是天降甘霖，是绝境中雪中送炭的大补药！

\vspace{2em}

惨剧，在瞬间发生。

惊恐的尖叫、绝望的哀嚎、徒劳的抵抗声，短暂地打破了山谷的死寂，又迅速归于更深的死寂。浓郁到化不开的血腥气冲天而起，却又被某种无形的力量贪婪地吞噬、吸收。片刻之后，原地只留下一地形容恐怖、血肉枯竭的干尸，他们携带的财物散落一地，在昏暗的天光下闪烁着冰冷而讽刺的光泽。

汲取了这数十人的精血，血魔苍白如纸的脸上恢复了一丝诡异的红润，周身溃散的气息也略微稳固。

他舔了舔嘴角并不存在的血迹，眼中红光闪烁。“荒域这群生灵，虽然品质低劣，但……足够用了。”

他辨认了一下方向，朝着灵气更稀薄、但在他看来“生机”可能更密集的东方，也就是大周国人口稠密的区域，继续他的亡命之旅。

\vspace{2em}

一个月，对于辽阔的大周国边境而言，是短暂而血腥的一个月。

先是边境附近的几个偏僻山村，一夜之间鸡犬不留，所有村民皆化为干尸，场面骇人听闻。地方官府惊惧，上报州府，派出的捕快和少量驻军进入山区查探，却如同泥牛入海，再无音讯。

紧接着，规模稍大的镇甸也遭了殃。防御的土墙木栅在血魔面前形同虚设，驻守的百十名兵丁连同镇民，尽数殒命。消息如同瘟疫般传开，边境线上人心惶惶，流言四起，有说山中出了专吸人血的妖魔，有说是敌国派来的诡秘军队。

朝廷震动，边关告急！地方驻军组织了一次规模较大的围剿，由一名经验丰富的将领率领上千精兵，携带强弓劲弩，开进事发区域。

然而，血魔虽伤势未愈，实力远未恢复巅峰，但其身为真正修士的手段，又岂是凡人军队所能理解？一场不对等的屠杀在夜幕下展开，军阵被诡异的血雾和鬼魅般的身法轻易撕破，将领阵亡，大军溃散，幸存者十不存一，带回了更加令人绝望和恐惧的消息。

妖魔！真正的、刀枪难伤、行动如风、嗜血如命的妖魔！

\vspace{2em}

朝堂之上，再次被沉重的阴云笼罩。官员们吵吵嚷嚷，却拿不出任何切实可行的办法。

调集更多军队？连上千精兵都全军覆没，再多的普通士卒上去恐怕也只是送死，还会导致边防空虚，给虎视眈眈的邻国可乘之机。

悬赏江湖异士？寻常武林高手，在那位仙师的徒弟面前都如同土鸡瓦狗，又能对这等妖魔有何作用？

白隐坐在龙椅上，听着下方一片慌乱无措的议论，眉头紧锁，指尖冰凉。他不由地又想起了那个总是神色平静的青年，想起了金殿柱上那道深达三寸的刻痕。或许……只有超越凡俗的力量，才能对付这种超越凡俗的灾厄？

下朝后，白隐独自在养心殿沉思良久，终于下定决心。他再次召来了白言冰。

“言冰，”白隐神色凝重，“边关妖魔之事，想必你已听闻。朝廷……已束手无策。朕知你与千羽跟随叶先生修行，已非凡俗。此事，关乎千万黎民性命，关乎国本安稳。朕……想请叶先生出手。”

白言冰早已从陈千羽那里得知了边境惨状，胸中亦是义愤填膺，闻言立刻道：“父皇放心，斩妖除魔，护佑百姓，本就是我辈应为！儿臣这就去禀明师尊！”

“不，”白隐抬手止住他，眼中闪过复杂神色，“这一次，朕需以国礼相请，方显郑重，亦绝他人非议。”他沉吟片刻，提笔写下一道敕令，加盖玉玺，“朕敕封叶牧谦先生为‘护国天师’，享国师礼遇，恳请天师出手，为民除害！”

当白言冰带着边境的详细情报提前一步回到小院时，叶牧谦正在指点陈千羽如何进一步提纯她那“淀粉丹”中的药力，去除杂质。

听完白言冰的叙述，尤其是听到“妖魔”、“吸食人血化为干尸”、“军队覆灭”等关键词时，叶牧谦的眼神陡然变得锐利起来。他放下了手中那颗卖相不佳的丹丸。

“修士……”他低声自语，心脏微微加速跳动。

这个词，终于以如此鲜明而残酷的方式，撞入了他的视野。

不同于之前模糊的猜测，这是活生生的、正在制造灾难的修士！这印证了他之前的猜想，我胡编的东西能被这两个徒弟走通，肯定是因为这方天地有猫腻！

但是新的问题是，其力量体系似乎与“问道途”迥异。

危险，同样是机遇。这或许是深入了解这个世界另一面的窗口，是检验“问道途”与真正修士力量对比的试金石。

当然，前提是，他们能对付得了。

随后，圣旨到了。叶牧谦自然是不跪的，仅仅是站了起来表示尊重。而传圣旨的天使也自然不敢有什么微词，于是便自顾自开念。

念完，叶牧谦也只是淡淡一笑。

“护国天师？虚名罢了。不过，此事我已知晓。妖魔肆虐，残害生灵，确不能坐视。”

他看向侍立一旁、眼中已燃起战意的白言冰，和面露忧色却同样坚定的陈千羽：“收拾一下，我们即刻动身，全速赶往边境。”

他没有拒绝皇帝的请求，甚至没有多提条件。于公，此事关乎无数百姓性命；于私，这是他必须面对的、来自世界另一面的第一次正式接触。他需要情报，需要验证。

白言冰精神一振，陈千羽也立刻点头，转身便去准备必要的丹药、干粮和随身之物。

次日清晨，三人动身，日夜兼程，穿州过府。

越靠近边境，空气中弥漫的那股若有若无的淡淡血腥与不祥气息便越发清晰。

沿途所见，村镇萧条，百姓面有菜色，眼神惊惶，偶尔看到官道上有零星逃难的人群，拖家带口，仓皇东去。白言冰面色沉凝，陈千羽眉宇间凝结着忧色，唯有叶牧谦，眼神深处除了凝重，还涌动着一丝难以言喻的探究。

根据最后的情报和残留的气息追踪，他们最终抵达了十万群山最东缘的一片荒凉山谷。这里曾是最后一个遭袭的边军营地旧址，焦黑的木桩、破碎的兵刃、以及早已干涸发黑、浸入泥土的大片污渍，无声地诉说着月前那场绝望的屠杀。

山谷中弥漫的死寂，比喧嚣更令人心悸。

“来了。”叶牧谦忽然停下脚步，目光投向山谷深处一片被阴影笼罩的乱石堆。

几乎在他出声的同时，乱石堆后，一道血红的身影如同鬼魅般飘然而出。

那是一个形貌极为诡异的男子。身材干瘦，面色是一种失血过多的惨白，偏偏双颊又泛着两团不正常的、仿佛涂抹上去的猩红。身上穿着一件破损不堪的暗红色袍子，上面沾满了深褐色的污渍，散发出浓烈的血腥与腐败混合的气味。

最引人注目的是他的眼睛，瞳孔竟是细长的竖瞳，泛着暗红色的幽光，如同盯上猎物的毒蛇。他周身缭绕着一层极淡的、令人极不舒服的暗红色雾气，仿佛有生命般微微蠕动。

正是肆虐边境月余的邪修——血魔。

血魔甫一现身，那双诡异的竖瞳便飞快地在叶牧谦三人身上扫过。他的目光首先落在为首的叶牧谦身上，略一感知，眼中便流露出一丝毫不掩饰的轻蔑与困惑。没有灵力波动，脊柱附近那至关重要的“灵枢”更是毫无感应，与这荒域的凡人一般无二。

随即他又看向白言冰和陈千羽，同样未能感知到灵枢的存在，只是觉得这两个年轻人气血似乎比寻常凡人旺盛一些，体魄也强健些，但依旧属于凡人范畴。

“呵……又是三个不知死活的凡人？是来追查同伙，还是来送饭的？”血魔的声音嘶哑干涩，如同砂纸摩擦，带着一种猫捉老鼠般的戏谑。他重伤未愈，虽靠吸食精血勉强恢复了些许，但远未到巅峰，此刻只想快点解决这几个“误入”的凡人，或许能再补充一点微不足道的消耗。

然而，站在他对面的白言冰，在血魔出现的刹那，全身肌肉便已悄然绷紧，进入了战斗状态。他养气初成，灵觉大幅提升，虽不懂什么真正的神识探查手段，但却能清晰地感觉到对方身上的异常。

没有气海！

师尊所传道途，迈入养气之后，力量核心便与丹田之下的气海息息相关。可眼前这个邪异的家伙，丹田之下空空如也，一片混沌。

但与此同时，白言冰却敏锐地察觉到，对方的脊柱部位，以及四肢的主要经络路径上，正流动着一股截然不同的、阴冷、暴戾且充满掠夺性的能量波动！那能量仿佛与他的血肉紧密相连，却又迥异于自身温养出的温润元气。这种矛盾而陌生的感知，让白言冰瞬间警惕到了极点。

“千羽，小心。此人……很怪。”白言冰低声对身旁已握紧一对钢针的陈千羽说道。

陈千羽微微点头，她魂骨已成，感知亦不弱，同样感觉到了那股令人极为不适的阴冷气息，仿佛能吸走生机。

“怪？”血魔听力极佳，嗤笑一声，“将死之人，废话倒多。”

他伤势虽然恢复大半，但仍然不欲久缠。身形一晃，竟化作一道模糊的血影，直扑看起来最弱的陈千羽！速度之快，远超寻常武林高手，同时一只干瘦如鸟爪的手掌探出，五指指尖泛起暗红邪光，带着腥风抓向陈千羽咽喉，意图一击毙命，吸摄精血。

“休想！”

白言冰厉喝一声，早已蓄势待发的他后发先至，脚步一错便挡在陈千羽身前，并指如剑，体内气海元气奔涌，循天骨滋生的内劲路径灌注于指尖，一记直刺，点向血魔抓来的手腕。这一指看似简单，却蕴含着养气境初成的沛然之力与天骨的沉凝穿透，指尖前方的空气发出轻微的呜咽。

血魔原本漫不经心，打算随手震开这“凡人”的指头，再抓碎那女子的喉咙。然而，当白言冰的指风及体时，他脸色骤变！

指尖传来的竟然是一股精纯、凝练、充满勃勃生机却又锐利无匹的奇异力量！这力量……甚至不是修士的灵力！

更让他心惊的是，对方攻击中蕴含的那股“意”，堂堂正正，隐隐对他周身的邪气产生了一种莫名的压制感！

“什么？！”血魔惊疑不定，抓出的手掌不得不变招，化爪为掌，暗红邪光暴涨，硬接了这一指。

“嗤——！”

指尖与掌心相触，发出一声如同烧红烙铁放入水中的嗤响。白言冰指间元气与血魔掌心血色邪气激烈碰撞、湮灭。

白言冰只觉一股阴寒刺骨、充满侵蚀性的力量顺着手臂经络试图钻入，他闷哼一声，气海元气自动流转，将那股侵入的邪气逼出、化去，但整条手臂仍是一阵酸麻。而血魔则感觉掌心传来灼痛，那凝聚的邪气竟被对方一指刺散大半，震得他手掌发麻，身形向后飘退两步，眼中充满了难以置信。

“你……你不是凡人！你这是什么力量？荒域怎么可能有修士？！”血魔失声叫道，竖瞳中第一次露出了惊骇。他固有的认知受到了剧烈冲击。

“荒域无修士”是内域的常识，这里灵气稀薄近乎于无，根本无法支撑灵枢径的修炼！可眼前这人，明明没有灵枢，没有灵力波动，却拥有着不弱于初入玄脉境修士的奇异力量！这完全颠覆了他的理解！

回答他的是陈千羽悄无声息袭来的钢针。她虽主修医道，但魂骨境界带来的身体掌控力与速度远超常人，抓住血魔惊疑后退的瞬间，钢针如毒蛇吐信，疾点血魔肋下、腰眼等要害。

招式简洁狠辣，专攻人体薄弱之处，同时她将一丝自身温养的、偏向生机的气机附着其上，试图干扰对方那阴邪的能量运行。

但血魔毕竟是真真正正的修士。血魔只守不攻、看似暂落下风，实际上却在暗中观察两人的招法，评估其实力。虽重伤未愈且惊骇，但血魔辗转腾挪的战斗本能仍在：身形诡异地一扭，如同没有骨头般，险险避开针刺，反手一掌拍向陈千羽肩头，掌风腥臭，邪气森然。

“你敢！”

陈千羽乃是白言冰的逆鳞，他的长剑终于出鞘！

一声清越龙吟，剑光如匹练般横斩而至，截向血魔手腕。这一剑，蕴含着他养气境的元气与天骨的沉猛力道，凌厉的剑气迫得血魔皮肤生疼。

血魔不得不再次撤掌后退，心中惊怒交加。

这两个年轻人，一个力量精纯沉凝，攻势堂堂正正却威力不俗；一个身法灵动，攻击刁钻且带着一股令他邪气运转微微滞涩的生机之力。两人配合默契，一时竟让他有些手忙脚乱；更让他憋屈的是，他完全看不懂对方的力量根源！

——但很快，血魔便断定，两人的实力一个顶多是玄脉境开脉期左右，一个估计只有灵枢境凝枢期巅峰，境界远低于他。虽然在如今情况之下，如果不全力以赴，别说恢复伤势，恐怕今日真要在这荒域阴沟里翻船了；但血魔也有信心，这两个家伙一定扛不住他全力施为！

“该死！不管你们是什么怪物，但区区玄脉境开脉期小修，还想阻拦我？今日，你们都要成为本座的血食！”

血魔彻底收起了轻视，眼中凶光爆射；低吼一声，周身暗红色邪气猛然暴涨，如同燃烧的血焰。随即，双手快速掐动一个诡异印诀，口中念念有词。

山谷中残留的淡淡血腥气仿佛受到了召唤，丝丝缕缕汇聚而来，融入他周身的血焰之中，使其威势再增；他本就惨白的脸色更白了几分，显然动用此法，对他负担也不小。

“血煞掌！”

血魔双掌齐出，两只由浓郁邪气凝聚而成的暗红色巨大掌印，带着刺鼻的血腥味和鬼哭般的尖啸，分别轰向白言冰和陈千羽！掌印所过之处，地面上的枯草迅速枯萎腐败，连岩石表面都留下了淡淡的腐蚀痕迹。

白言冰面色凝重，他能感受到血魔认真了，这一击的威力远超之前。他深吸一口气，气海元气奔腾，尽数灌注于长剑之上，剑身嗡鸣，泛起一层淡金色的微光。他迎着一道掌印，将全身力量、意志与对守护之念的感悟融于一剑，笔直刺出！

“破！”

剑尖与血色掌印狠狠对撞！

“轰——！”

一声闷响，气劲四溢。白言冰浑身剧震，虎口迸裂，鲜血染红剑柄，身形踉跄后退数步，体内气血翻腾，元气一阵紊乱。那血色掌印也被这一剑刺穿、绞碎大半，但残余的邪气依然冲击得他颇为难受。

另一边，陈千羽面对另一道掌印，自知硬接不下，将魂骨境界的身法发挥到极致，如同风中柳絮般飘忽闪避，同时不断调动暴雨梨花般的钢针，削弱掌印边缘的邪气。

但掌印范围太大，邪气又极具侵蚀性，几缕漏过的邪气擦过她的手臂，顿时传来火辣辣的刺痛，仿佛生机被强行抽走一丝，手臂肌肤也瞬间变得灰暗了一些。她闷哼一声，脸色发白。

仅仅一击，高下立判。血魔全力施为之下，即便重伤，玄脉境的威力展露无遗，白言冰二人顿时落入下风，险象环生。

血魔得势不饶人。看起来他的估计并没有错：拼上限，还是他更高一筹！

狞笑着，血魔身形再动，化作道道血影，从四面八方攻向两人。

邪气纵横，爪风凌厉，招招夺命。

白言冰奋力挥剑格挡，剑光与血影不断碰撞，爆出一团团气劲，但他左支右绌，身上很快多了几道深可见骨的伤口。伤口处血肉翻卷，却诡异地没有多少鲜血流出，反而呈现出灰败之色，是邪气不断侵蚀的结果。

陈千羽境界更低，也是只能凭借身法狼狈躲避，偶尔反击也被轻易化解，险象环生。若非白言冰多次拼死相救，早已香消玉殒。

“师尊！”白言冰又一次格开血魔抓向陈千羽头颅的一爪，自己后背却被邪气扫中，喷出一口鲜血，嘶声喊道。他知道，仅凭他们二人，绝非这邪修对手。

一直在后方静静观战、仿佛置身事外的叶牧谦，此刻终于动了。

他一直在观察，观察血魔的力量运行方式，观察这群修士战斗的特点，也与自身问道途暗暗对比；更是让他们过一次生死门，在实战中强化修为与感悟，也是告诉他们修仙界并不是什么良善之地。

而此刻，见徒弟们确实已到极限，是该出手了。

他向前踏出一步，步伐不大，却仿佛瞬间缩短了与战场的距离。面对狂扑而来、意图先解决白言冰的血魔，叶牧谦只是平平无奇地抬起右手，轻飘飘一掌推出。

这一掌，没有惊天动地的声势，没有光华耀眼的异象，甚至没有带起多大的掌风。叶牧谦谨记反噬教训，只动用了大约相当于自身此刻气海容量一成左右的力量——当然，这一成力已经比白言冰的全部还多了。

就在掌力发出的瞬间，异变陡生！

叶牧谦只觉自己冥冥中仿佛触及了某种更深层的东西。自己那基于内求、温养而生的元气，在离体之后，竟隐隐引动了周遭天地间一丝极其微弱的、难以言喻的共鸣？或者说，他这问道途的力量，似乎对灵枢径这种明显偏向掠夺、借外的力量，有着某种先天性的压制？

血魔原本对叶牧谦这看似毫无威胁的一掌不屑一顾，正打算随手拍散，然后继续虐杀那两个让他吃惊的小辈。但当那掌力及身时，他脸上的狞笑瞬间凝固，化为无边的恐惧！

那掌力之中，蕴含着一股他无法理解、无法抗拒的意志！

那是一种层次上的绝对差距，仿佛蝼蚁面对苍穹，邪祟直面天威！

他周身沸腾的血煞邪气，在这掌力面前，竟如同雪遇骄阳，迅速消融、溃散！

他赖以生存、催动的灵枢与玄脉中的邪力，运转瞬间凝滞，仿佛被冻结！

“不……不可能！怎么可能会有真符境！”血魔惊恐万状地尖叫，拼命运转所有邪力，在身前布下一层又一层暗红色的邪气护盾。

在他的认知中，能轻易压制于他的人似乎也只有所谓的“真符境”了。

“噗、噗、噗……”

轻响连连，那层层邪气护盾在那平淡无奇的掌力面前，比纸糊的还要脆弱，一触即溃。

掌力毫无阻碍地印在了血魔仓促架起的双臂，以及他的胸膛之上。

没有血肉横飞，只有“砰”的一声巨响。

血魔整个人随即倒飞出去，双眼瞬间失去神采。他整个人的动作彻底僵住，周身的邪气如同潮水般退去、消散。紧接着，他干瘦的身体如同风化的沙雕，从被掌力击中的部位开始，迅速蔓延出无数细密的裂纹。

“喀嚓……哗啦……”

裂纹遍布全身，下一刻，血魔整个人竟化作一蓬暗红色的细灰，簌簌洒落在地，连半点残骸都未曾留下。唯有原地留下一个浅浅的掌印凹坑，以及几件物品：一个巴掌大小、绣着诡异血色纹路的灰色布袋，一块刻着“血魔”二字的暗红色金属令牌，以及几件零碎的、散发着微弱邪气的衣服。

\vspace{2em}

山谷中，死一般的寂静。

白言冰柱剑而立，喘着粗气，身上伤口还在渗着诡异的灰败液体。陈千羽捂着受伤的手臂，脸色苍白。两人都瞪大了眼睛，难以置信地看着眼前这一幕。

就这么……结束了？

那让他们拼尽全力、险死还生都无法抗衡的恐怖邪修，在师尊那轻描淡写的一掌之下，竟然……灰飞烟灭？

叶牧谦自己也有些愕然。他看了看自己的手掌，又看了看地上那滩灰烬和掌印。一成力？他本以为最多重伤对方，没想到直接拍成了灰……看来，这元气的力量性质，比自己想象的还要特殊，对这类邪修似乎有着超乎想象的克制？

或者说……是那冥冥中引动的“天威”般的意蕴在起作用？

白言冰喉结滚动了几下，终于问道：“师尊，你用了多大力道？”

叶牧谦回过神来：“一成……？”

寂静，死一样的寂静。

\vspace{2em}

叶牧谦摇了摇头，走到那几件遗物前。先捡起那块令牌，触手冰凉，材质不知是什么，正面是“血魔”两个古篆，背面则是一些扭曲难明的符文，依稀辨认是“血魂教副执事”六个字，散发着淡淡的邪气与血腥味。

他又拿起那个灰色布袋，入手轻若无物，但质地坚韧，尝试将一丝元气探入，竟感觉里面有一个大约三尺见方的奇异空间！

里面散落着一些瓶瓶罐罐、几叠符纸、几块暗红色的晶石，约有几两之数的银锭，以及一些杂七杂八的东西。这些东西，无一例外都缭绕着或浓或淡的邪恶气息，令人望之生厌。

除此之外，还有些似玉的美石，其上邪恶气息并不充盈。

在叶牧谦猜测中，这大概就是前世许多玄幻小说中，修士们的通用货币——灵石？

“师尊，这些……”白言冰走上前，看着那些邪气森然的物品，眉头紧皱。

“邪修之物，大多以生灵精血、魂魄等歹毒材料炼制，有伤天和。”叶牧谦淡淡道，他虽不知具体炼制法门，但那股令人极度不适的邪恶气息做不得假。

他看向陈千羽，“千羽，你觉得呢？”

陈千羽强忍着伤口处邪气侵蚀带来的麻木与刺痛，仔细感知了一下那些丹药和符箓的气息，作为医师，她对“生机”与“死气”、“正气”与“邪气”的感知尤为敏锐。

她脸色难看地摇头：“这些丹药，蕴含暴烈死气与怨念，服之恐瞬间侵蚀生机，沦为只知杀戮的怪物。这些符箓，亦是引动血煞、污秽之力的邪物。留之反是祸患。”

叶牧谦点头：“既如此，毁了吧。”

他倒空储物袋，收起那些灵石和银两。随即运起一丝元气，包裹住那些邪丹、邪符以及暗红晶石，轻轻一催。

只听一阵轻微的“嗤嗤”声，那些邪物上的气息迅速消散，丹药化为黑粉，符箓自燃成灰，晶石也碎裂开来，失去光泽。

唯有那个材质特殊的储物袋和那枚“血魔”令牌留存下来。令牌或许可作为击杀凭证，有些修士通行的货币总是好的，储物袋则无疑是个实用的好东西，虽然沾染过邪气，但以元气反复冲刷后，应可勉强使用。

处理完这些，叶牧谦看向陈千羽：“你的伤如何？”

陈千羽正以银针刺穴，辅以自身温养的生机之气，努力逼出、化解侵入伤口的邪气。闻言，她微微蹙眉：“邪气阴毒，侵蚀生机，好在量不多，弟子可以处理。”

在逼出邪气、修复伤口的过程中，她对抗着那股掠夺生机的力量，对生机的流转、平衡、以及被强行破坏后的反应，有了比以往任何一次医治都更深刻、更直接的体会。生与死，正与邪，在这种对抗中显得如此鲜明。

不知不觉间，她体内那魂骨的瓶颈悄然松动，骨骼深处那温润的光泽变得更加凝实、深邃，隐隐向着更坚韧、更通透的层次转化——正是天骨境界的特征！战斗与疗伤的双重感悟，竟让她在此刻水到渠成地突破了。

白言冰也察觉到了妹妹气息的变化，惊喜道：“千羽，你突破了？”

陈千羽缓缓吐出一口浊气，点了点头，脸上露出一丝疲惫却欣慰的笑容：“嗯，侥幸。”

\vspace{2em}

喜悦过后，更大的疑惑却涌上师徒三人心头。

白言冰看着地上那滩灰烬，沉声道：“师尊，此人……很是古怪。弟子完全感觉不到他有气海的存在，但他的脊柱和四肢主要经络中，却流动着那种阴邪的能量。难道邪修的修炼方式，都是如此迥异吗？”

陈千羽也补充道：“而且他惊呼‘荒域无修士’，称我们为‘怪物’，还提到了‘灵枢’……师尊，他口中的‘荒域’、‘内域’、‘灵枢径’究竟是何意？”

叶牧谦沉默了片刻，看着两双充满求知与困惑的眼睛，最终缓缓摇了摇头，坦诚道：“不知道。”

他是真的不知道。血魔死得太快，没能拷问出任何情报。这些名词，他也是第一次从这邪修口中听说。

“但是，”叶牧谦话锋一转，目光变得悠远，望向十万群山那苍茫的轮廓，“有一点可以确定。在群山之外，那个被称为‘内域’的地方，确实存在着一个修士的社会，他们修炼的，很可能就是这种以‘灵枢’为核心、与我们道途截然不同的‘灵枢径’。此人，不过是其中修炼邪法、被追杀逃窜至此的一个。”

他顿了顿，看着自己的两个徒弟，语气郑重：“人外有人，天外有天。我们之前所见，不过一隅。真正的天地，远比我们想象的广阔，也藏着我们无法理解的奥秘与危险。”

白言冰和陈千羽对视一眼，都从对方眼中看到了震撼、了然，以及一种被点燃的、熊熊燃烧的探索之火。京城朝堂的倾轧、边境百姓的苦难、甚至大周国的疆域，在这突如其来的、来自世界另一面的信息冲击下，似乎都变得渺小了。

“师尊，”白言冰握紧了手中的剑，伤口虽疼，眼神却无比明亮，“我们……要出去看看吗？去看看那‘内域’，去看看那‘灵枢径’，去看看这真正的天地？”

陈千羽虽未说话，但那双清亮的眸子中，同样写满了坚定与向往。

叶牧谦看着他们，嘴角微微勾起一丝弧度。他来到这个世界，懵懂间开创了“问道途”，收了两个徒弟，经历了诸多变故。原本或许只想偏安一隅，弄清自身奥秘。但此刻，一条更广阔、更波澜壮阔的道路，似乎已铺展在眼前。

而且他也确定呆在这里不能让他回家。

“当然。”他轻声答道，目光仿佛已穿越了重重山峦，“不过，在离开之前，我们还需回京复命，做些准备。而且……”

他看了一眼陈千羽手臂上尚未完全祛除邪气的伤口，以及白言冰身上多处灰败的创伤。

“先治好伤。前面的路，恐怕不会太平。”

山谷的风吹过，卷起一丝灰烬，也送来了远方群山更深处的、未知的气息。

\vspace{2em}

边境山谷一战，尘埃终于落定。叶牧谦师徒三人将邪修“血魔”伏诛的消息通过边军驿站急报京城后，又参与了些协助边军与难民打扫战场、安葬死者的工作，便踏上了归程。

归途不似来时紧迫，但三人心中却比来时更为沉重，也更为明亮。沉重的是亲眼所见的惨状与那迥异修炼体系带来的冲击，明亮的则是前方豁然开朗的、名为“内域”的全新天地。

数日后，京城在望。他们没有惊动太多人，径直入宫面圣。

养心殿内，白隐早已接到捷报，正翘首以盼。见到三人安然归来，尤其是看到白言冰身上虽已愈合但痕迹犹存的伤口，以及陈千羽眉宇间那丝经历生死搏杀后的坚毅，这位帝王心中百感交集，既有如释重负的欣慰，也有难以言喻的酸涩。

叶牧谦简略汇报了诛杀“血魔”的经过，并呈上了那枚“血魔”令牌作为凭证。白隐摩挲着那冰凉的令牌，听着那凶魔伏诛的过程，长长舒了一口气，仿佛卸下了心头一块巨石。边境无数百姓的性命，上千将士的亡魂，总算有了一个交代。

“叶先生，言冰，千羽……你们，又救了大周一次。”白隐的声音有些沙哑，目光复杂地扫过三人，“朕……不知该如何赏赐你们了。”

叶牧谦微微摇头：“陛下，分内之事，无需赏赐。倒是有一事，需向陛下禀明。”

他看了一眼身旁并肩而立的白言冰和陈千羽，两人会意，上前一步，齐齐跪下。

白言冰沉声道：“父皇，边境妖魔虽除，但此事让儿臣与千羽明白，天地之广阔，远超我等想象。十万群山之外，另有世界，其中修士纵横，力量体系与我等所学截然不同。儿臣与千羽心意已决，欲随师尊跨越群山，前往那‘内域’游历修行，探寻大道，也为将来可能再临的威胁，积累见识与力量。今日……特来向父皇辞行。”

陈千羽亦叩首，声音轻柔却坚定：“陛下多年养育、庇护之恩，千羽永世不忘。然师尊所言‘人外有人，天外有天’，千羽深以为然。医道无穷，弟子欲往更广阔的天地求索，亦想看看那不同的修炼文明，于生死、正邪之道上有何见解。请陛下……允准。”

白隐愣住了。他虽知儿女跟随叶牧谦修行，迟早会离开，却没想到这一天来得如此之快，且目标竟是那凶险莫测、刚刚夺走无数人性命的群山之外。

他张了张嘴，想说什么，目光触及叶牧谦平静深邃的眼神，又看向跪在下方、眼神中充满不容动摇的决心的儿子和（养）女，千言万语堵在喉头，最终化作一声悠长的叹息。

他起身，绕过御案，亲手将两人扶起。他的手有些颤抖，仔细端详着白言冰愈发刚毅的面庞和陈千羽清丽却坚执的眉眼，仿佛要将他们的样子刻在心里。

“你们的未来……确实不在朝廷啊。”白隐喃喃道，眼中有不舍，有担忧，但最终被一种更为复杂的情绪取代——那是看到雏鹰终于要振翅翱翔于真正苍穹时的释然与骄傲。

“跟着叶先生，海阔天空，去追寻你们自己的道吧。朕……为父的若是继续干涉，也只会成为你们的拖累。”

他拍了拍白言冰坚实的肩膀，又轻轻抚了抚陈千羽的发顶，如同寻常人家的老父亲。“去吧。京城，大周，有为父在。你们只需记得，无论走多远，这里永远是你们的根。若累了，想回来了，家门永远为你们敞开。”

他又转向叶牧谦，郑重地躬身一礼：“叶先生，朕这一双儿女，就托付给您了。他们性子执拗，若有不当之处，万望先生海涵，多加指点。”

叶牧谦侧身避过，颔首道：“陛下言重了。言冰与千羽皆是我弟子，我自当尽心。”

辞别之事既定，便不再拖泥带水。

白隐执意要设家宴饯行，席间只有帝后、太子白承烨与师徒三人。没有朝堂的规矩，只有家人间的絮语与叮嘱。

次日，天光未亮，三人便已收拾妥当。行装极其简单，除了必要的衣物、干粮、清水，便是陈千羽精心准备的各类丹药，白言冰的长剑，以及那个从血魔处得来、已被叶牧谦以元气反复冲刷过的储物袋。他们谢绝了白隐派兵护送的好意，如同寻常旅人，悄然从侧门离开了这座居住多年的巍峨皇城。

回首望去，晨曦中的宫阙轮廓渐渐模糊。白言冰与陈千羽驻足片刻，眼中皆有感慨，但更多的是对前路的坚定。叶牧谦亦看了一眼那象征着世俗权力巅峰的建筑群，心中波澜微起，随即平复。这里的故事告一段落，新的篇章，将在山的那一边展开。

三人离了京城，一路向西，朝着十万群山的方向行去。起初尚有官道驿站，越往西行，人烟越见稀少，地势渐趋崎岖，景致也从繁华城镇变为旷野疏林，最终踏入连绵的丘陵地带。空气似乎也清新冷冽了许多，带着山野特有的气息。

“看，师尊，绝念峰。”白言冰忽然指着一座视野内的山峰道。

翻过绝念，便再难回。

“嗯。”叶牧谦的脸上看不出心理变化。

这是他来到此世的第一站。而今天，他就要从这里离开荒域了。

也许人生就是这样，一节一节的过下去。

“休息片刻，继续赶路吧。”他的声音平静却带着期许，“前面的山，该翻了。”